%! TeX program = pdflatex
\documentclass[aspectratio=169,12pt]{beamer}

\usepackage[T1]{fontenc}
\usepackage{lmodern}
\usepackage[default]{sourcesanspro}
\usefonttheme[onlymath]{serif}
\usepackage{amsmath,amssymb,mathtools}
\usepackage{booktabs,array,microtype}
\usepackage{tikz}
\usetikzlibrary{
  arrows.meta,
  backgrounds,
  calc,
  decorations.pathreplacing,
  fit,
  positioning,
  shapes.geometric
}
\usepackage{appendixnumberbeamer}
\usepackage{beamerthemethesis}

\newcommand{\OPT}{\mathrm{OPT}}
\newcommand{\optpost}{\OPT_{\mathrm{post}}}
\newcommand{\optadapt}{\OPT_{\mathrm{adapt}}}
\newcommand{\optprior}{\OPT_{\mathrm{prior}}}
\newcommand{\E}{\mathbb{E}}
\newcommand{\Pp}{\mathbb{P}}
\newcommand{\cost}{\operatorname{cost}}
\newcommand{\poly}{\operatorname{poly}}

% A more saturated blue reserved for thesis contributions in comparison
% tables.  The regular theme blue remains unchanged elsewhere.
\definecolor{ThResultBlue}{HTML}{0057B8}

% The same information-gap table is reused as a progress marker throughout
% the talk.  Stage 0 shows only prior work; stages 1--4 cumulatively reveal
% the four new adaptivity/clairvoyance cells and highlight the newest one.
\newcommand{\thgapbox}[2][]{%
  \tikz[baseline=(thgapcell.base)]{%
    \node[
      minimum width=3.00cm,
      minimum height=0.88cm,
      inner xsep=3pt,
      inner ysep=2pt,
      align=center,
      rounded corners=1pt,
      #1
    ] (thgapcell) {#2};
  }%
}

\newcommand{\thgapstagecell}[3]{%
  \ifnum#1<#2\relax
    \thgapbox[opacity=0]{#3}% not introduced yet
  \else
    \ifnum#1=#2\relax
      \thgapbox[draw=ThResultBlue,fill=ThResultBlue!10,line width=1.4pt]{#3}% newly introduced
    \else
      \thgapbox{#3}% introduced on an earlier slide
    \fi
  \fi
}

\newcommand{\ThGapProgressTable}[1]{%
  \centering
  {\fontsize{10.8pt}{13.0pt}\selectfont
  \renewcommand{\arraystretch}{1.55}
  \setlength{\tabcolsep}{8pt}
  \begin{tabular}{>{\bfseries}lccc}
    \toprule
    Metric
      & \shortstack{Obliviousness gap\\[-0.04cm]
          \(\optprior/\optpost\)}
      & \shortstack{Adaptivity gap\\[-0.04cm]
          \(\optprior/\optadapt\)}
      & \shortstack{Clairvoyance gap\\[-0.04cm]
          \(\optadapt/\optpost\)} \\
    \midrule
    Symmetric
      & \(\le 3^{[1]}\)
      & \thgapstagecell{#1}{1}{%
          \shortstack{
            \textcolor{ThResultBlue}{\(\boldsymbol{\ge 2}\)}\\[-0.04cm]
            \(\le 3^{[1]}\)}}
      & \thgapstagecell{#1}{3}{%
          \shortstack{
            \textcolor{ThResultBlue}{\(\boldsymbol{\ge 10/9}\)}\\[-0.04cm]
            \(\le 3^{[1]}\)}} \\
    \cmidrule(lr){1-4}
    Asymmetric
      & \(\begin{gathered}
            \Omega(n^{1/4}/\log n)\\[-0.04cm]
            O(\sqrt n)
          \end{gathered}^{[2]}\)
      & \thgapstagecell{#1}{2}{%
          \shortstack{
            \textcolor{ThResultBlue}{\(\boldsymbol{\Omega(n^{1/3}/\log n)}\)}\\[-0.04cm]
            \(O(\sqrt n)^{[2]}\)}}
      & \thgapstagecell{#1}{4}{%
          \shortstack{
            \textcolor{ThResultBlue}{\(\boldsymbol{\ge 3/2}\)}\\[-0.04cm]
            \textcolor{ThResultBlue}{\(\boldsymbol{O(n^{1/3})}\)}}} \\
    \bottomrule
  \end{tabular}}

  \vfill
  \begin{minipage}{0.94\linewidth}
    \raggedright
    {\fontsize{7.5pt}{9pt}\selectfont\color{ThMuted}
      [1] Shmoys and Talwar (2008).\par
      [2] Christalla, Puhlmann, and Traub (2025).\par
      \ifnum#1>0\relax
        \textcolor{ThResultBlue}{\bfseries Bold blue}: this thesis.
        Regular black: previous work or an inherited upper bound.
      \else
        \phantom{Bold blue: this thesis.
        Regular black: previous work or an inherited upper bound.}
      \fi}
  \end{minipage}
}

\title[Adaptivity Bounds for Stochastic TSP]{%
  Adaptivity Bounds for the\\
  Stochastic Traveling Salesperson Problem}
\author{Teng Liu}
\institute{Department of Computer Science \enspace\textperiodcentered\enspace ETH Z\"urich}
\date{Master Thesis Defense \enspace\textperiodcentered\enspace 2026}

% Native TikZ diagrams for the thesis defense.

\tikzset{
  thDet/.style={
    rectangle,rounded corners=0.45pt,draw=ThInk,fill=ThWhite,
    line width=0.7pt,minimum width=0.52cm,minimum height=0.48cm,
    inner sep=1.5pt,font=\fontsize{8.5pt}{9pt}\selectfont
  },
  thDepot/.style={
    thDet,fill=ThGold!24,draw=ThGold!65!ThInk,line width=0.9pt
  },
  thStoch/.style={
    circle,draw=ThInk,fill=ThTeal!18,line width=0.7pt,
    minimum size=0.47cm,inner sep=1pt,
    font=\fontsize{8pt}{9pt}\selectfont
  },
  thActive/.style={
    thStoch,fill=ThTeal,text=ThWhite,draw=ThTeal!55!ThInk
  },
  thPacket/.style={
    thStoch,fill=ThGold,text=ThWhite,draw=ThGold!55!ThInk,
    line width=0.85pt
  },
  thPacketActive/.style={
    thPacket,draw=ThTeal,line width=1.25pt
  },
  thInactive/.style={
    thStoch,fill=ThPaper,draw=ThMuted,text=ThMuted
  },
  thArrow/.style={
    -{Stealth[length=2.2mm,width=1.45mm]},line width=1.05pt,draw=ThBlue
  },
  thArrowThin/.style={
    -{Stealth[length=1.8mm,width=1.2mm]},line width=0.72pt,draw=ThBlue
  },
  thEdge/.style={draw=ThInkSoft,line width=0.72pt},
  thFaint/.style={draw=ThLine,line width=0.55pt},
  thDashed/.style={draw=ThMuted,densely dashed,line width=0.72pt},
  thLabel/.style={
    font=\fontsize{8.5pt}{10pt}\selectfont,text=ThMuted,align=center
  },
  thSmallLabel/.style={
    font=\fontsize{7.5pt}{9pt}\selectfont,text=ThMuted,align=center
  }
}

% --------------------------------------------------------------------------
% Classical metric TSP: a complete graph with one tour highlighted.
% --------------------------------------------------------------------------
\newcommand{\ThTSPGraph}{%
\begin{tikzpicture}[x=1cm,y=1cm]
  \coordinate (r) at (-2.15,0.00);
  \coordinate (a) at (-1.10,1.48);
  \coordinate (b) at (0.72,1.42);
  \coordinate (c) at (2.02,0.20);
  \coordinate (d) at (0.92,-1.42);
  \coordinate (e) at (-1.02,-1.50);

  \foreach \u/\v in {r/a,r/b,r/c,r/d,r/e,a/b,a/c,a/d,a/e,b/c,b/d,b/e,c/d,c/e,d/e}{
    \draw[ThMuted!42,line width=0.55pt] (\u)--(\v);
  }

  \draw[ThBlue,line width=2.05pt,line join=round]
    (r)--(a)--(b)--(c)--(d)--(e)--cycle;

  \node[circle,draw=ThGold!72!ThInk,fill=ThGold!22,
    line width=0.9pt,minimum size=0.48cm,inner sep=0pt,
    font=\fontsize{9pt}{10pt}\selectfont] at (r) {$r$};
  \foreach \v in {a,b,c,d,e}{
    \node[circle,draw=ThInkSoft,fill=ThWhite,line width=0.8pt,
      minimum size=0.36cm,inner sep=0pt] at (\v) {};
  }
\end{tikzpicture}}

% --------------------------------------------------------------------------
% Stochastic TSP: four click-through activation realizations.
% --------------------------------------------------------------------------
\newcommand{\ThStochasticGraph}{%
\begin{tikzpicture}[x=1cm,y=1cm]
  \path[use as bounding box] (-2.45,-1.82) rectangle (2.30,1.82);
  \coordinate (r) at (-2.15,0.00);
  \coordinate (a) at (-1.10,1.48);
  \coordinate (b) at (0.72,1.42);
  \coordinate (c) at (2.02,0.20);
  \coordinate (d) at (0.92,-1.42);
  \coordinate (e) at (-1.02,-1.50);

  \foreach \u/\v in {r/a,r/b,r/c,r/d,r/e,a/b,a/c,a/d,a/e,b/c,b/d,b/e,c/d,c/e,d/e}{
    \draw[ThMuted!34,line width=0.55pt] (\u)--(\v);
  }

  \node[circle,draw=ThGold!72!ThInk,fill=ThGold!22,
    line width=0.9pt,minimum size=0.50cm,inner sep=0pt,
    font=\fontsize{9pt}{10pt}\selectfont] at (r) {$r$};

  \only<1>{
    \node[thActive,minimum size=0.50cm] at (a) {$1$};
    \node[thInactive,minimum size=0.50cm] at (b) {$2$};
    \node[thActive,minimum size=0.50cm] at (c) {$3$};
    \node[thInactive,minimum size=0.50cm] at (d) {$4$};
    \node[thActive,minimum size=0.50cm] at (e) {$5$};
  }
  \only<2>{
    \node[thInactive,minimum size=0.50cm] at (a) {$1$};
    \node[thActive,minimum size=0.50cm] at (b) {$2$};
    \node[thActive,minimum size=0.50cm] at (c) {$3$};
    \node[thActive,minimum size=0.50cm] at (d) {$4$};
    \node[thInactive,minimum size=0.50cm] at (e) {$5$};
  }
  \only<3>{
    \node[thActive,minimum size=0.50cm] at (a) {$1$};
    \node[thInactive,minimum size=0.50cm] at (b) {$2$};
    \node[thInactive,minimum size=0.50cm] at (c) {$3$};
    \node[thActive,minimum size=0.50cm] at (d) {$4$};
    \node[thActive,minimum size=0.50cm] at (e) {$5$};
  }
  \only<4>{
    \node[thActive,minimum size=0.50cm] at (a) {$1$};
    \node[thActive,minimum size=0.50cm] at (b) {$2$};
    \node[thInactive,minimum size=0.50cm] at (c) {$3$};
    \node[thInactive,minimum size=0.50cm] at (d) {$4$};
    \node[thActive,minimum size=0.50cm] at (e) {$5$};
  }
\end{tikzpicture}}

% --------------------------------------------------------------------------
% A posteriori: reveal the complete active set, then optimize for it.
% --------------------------------------------------------------------------
\newcommand{\ThPosterioriExample}{%
\begin{tikzpicture}[x=1cm,y=1cm]
  \path[use as bounding box] (0,-1.45) rectangle (11,1.45);
  \begin{scope}[shift={(2.35,0.20)},scale=0.70]
    \coordinate (lr) at (-2.15,0.00);
    \coordinate (lone) at (-1.10,1.48);
    \coordinate (ltwo) at (0.72,1.42);
    \coordinate (lthree) at (2.02,0.20);
    \coordinate (lfour) at (0.92,-1.42);
    \coordinate (lfive) at (-1.02,-1.50);

    \foreach \u/\v in {lr/lone,lr/ltwo,lr/lthree,lr/lfour,lr/lfive,
      lone/ltwo,lone/lthree,lone/lfour,lone/lfive,
      ltwo/lthree,ltwo/lfour,ltwo/lfive,
      lthree/lfour,lthree/lfive,lfour/lfive}{
      \draw[ThMuted!34,line width=0.55pt] (\u)--(\v);
    }

    \node[circle,draw=ThGold!72!ThInk,fill=ThGold!22,
      line width=0.9pt,minimum size=0.50cm,inner sep=0pt,
      font=\fontsize{9pt}{10pt}\selectfont] at (lr) {$r$};
    \node[thActive,minimum size=0.50cm] at (lone) {$1$};
    \node[thInactive,minimum size=0.50cm] at (ltwo) {$2$};
    \node[thActive,minimum size=0.50cm] at (lthree) {$3$};
    \node[thInactive,minimum size=0.50cm] at (lfour) {$4$};
    \node[thActive,minimum size=0.50cm] at (lfive) {$5$};
  \end{scope}
  \node[thLabel,text=ThTeal] at (2.35,-1.30) {$A=\{1,3,5\}$};

  \draw[thArrow,draw=ThTeal,line width=1.35pt]
    (4.65,0.22)--(6.15,0.22);

  \begin{scope}[shift={(8.55,0.20)},scale=0.70]
    \coordinate (rr) at (-2.15,0.00);
    \coordinate (rone) at (-1.10,1.48);
    \coordinate (rthree) at (2.02,0.20);
    \coordinate (rfive) at (-1.02,-1.50);

    \foreach \u/\v in {rr/rone,rr/rthree,rr/rfive,
      rone/rthree,rone/rfive,rthree/rfive}{
      \draw[ThMuted!30,line width=0.55pt] (\u)--(\v);
    }
    \draw[ThTeal,line width=2.05pt,line join=round]
      (rr)--(rone)--(rthree)--(rfive)--cycle;

    \node[circle,draw=ThGold!72!ThInk,fill=ThGold!22,
      line width=0.9pt,minimum size=0.50cm,inner sep=0pt,
      font=\fontsize{9pt}{10pt}\selectfont] at (rr) {$r$};
    \node[thActive,minimum size=0.50cm] at (rone) {$1$};
    \node[thActive,minimum size=0.50cm] at (rthree) {$3$};
    \node[thActive,minimum size=0.50cm] at (rfive) {$5$};
  \end{scope}
  \node[thLabel,text=ThTeal] at (8.55,-1.30) {cost $=\OPT(A)$};
\end{tikzpicture}}

% --------------------------------------------------------------------------
% A priori: fix a full order, then shortcut the inactive vertices.
% --------------------------------------------------------------------------
\newcommand{\ThPrioriExample}{%
\begin{tikzpicture}[x=1cm,y=1cm]
  \path[use as bounding box] (0,-1.45) rectangle (11,1.45);
  \begin{scope}[shift={(2.35,0.20)},scale=0.70]
    \coordinate (lr) at (-2.15,0.00);
    \coordinate (lone) at (-1.10,1.48);
    \coordinate (ltwo) at (0.72,1.42);
    \coordinate (lthree) at (2.02,0.20);
    \coordinate (lfour) at (0.92,-1.42);
    \coordinate (lfive) at (-1.02,-1.50);

    \foreach \u/\v in {lr/lone,lr/ltwo,lr/lthree,lr/lfour,lr/lfive,
      lone/ltwo,lone/lthree,lone/lfour,lone/lfive,
      ltwo/lthree,ltwo/lfour,ltwo/lfive,
      lthree/lfour,lthree/lfive,lfour/lfive}{
      \draw[ThMuted!30,line width=0.55pt] (\u)--(\v);
    }
    \draw[ThTeal,line width=2.05pt,line join=round]
      (lr)--(lone)--(ltwo)--(lthree)--(lfour)--(lfive)--cycle;

    \node[circle,draw=ThGold!72!ThInk,fill=ThGold!22,
      line width=0.9pt,minimum size=0.50cm,inner sep=0pt,
      font=\fontsize{9pt}{10pt}\selectfont] at (lr) {$r$};
    \foreach \pos/\lab in {lone/1,ltwo/2,lthree/3,lfour/4,lfive/5}{
      \node[circle,draw=ThInkSoft,fill=ThPaper,line width=0.8pt,
        minimum size=0.50cm,inner sep=0pt,
        font=\fontsize{9pt}{10pt}\selectfont] at (\pos) {$\lab$};
    }
  \end{scope}
  \node[thLabel,text=ThTeal] at (2.35,-1.30)
    {$T=(r,1,2,3,4,5,r)$};

  \draw[thArrow,draw=ThTeal,line width=1.35pt]
    (4.65,0.22)--(6.15,0.22);
  \node[thLabel,text=ThTeal,fill=ThPaper,inner sep=1.5pt]
    at (5.40,0.64) {$A=\{2,4,5\}$};

  \begin{scope}[shift={(8.55,0.20)},scale=0.70]
    \coordinate (rr) at (-2.15,0.00);
    \coordinate (rone) at (-1.10,1.48);
    \coordinate (rtwo) at (0.72,1.42);
    \coordinate (rthree) at (2.02,0.20);
    \coordinate (rfour) at (0.92,-1.42);
    \coordinate (rfive) at (-1.02,-1.50);

    \foreach \u/\v in {rr/rone,rr/rtwo,rr/rthree,rr/rfour,rr/rfive,
      rone/rtwo,rone/rthree,rone/rfour,rone/rfive,
      rtwo/rthree,rtwo/rfour,rtwo/rfive,
      rthree/rfour,rthree/rfive,rfour/rfive}{
      \draw[ThMuted!25,line width=0.55pt] (\u)--(\v);
    }
    \draw[ThTeal,line width=2.05pt,line join=round]
      (rr)--(rtwo)--(rfour)--(rfive)--cycle;

    \node[circle,draw=ThGold!72!ThInk,fill=ThGold!22,
      line width=0.9pt,minimum size=0.50cm,inner sep=0pt,
      font=\fontsize{9pt}{10pt}\selectfont] at (rr) {$r$};
    \node[thInactive,minimum size=0.50cm] at (rone) {$1$};
    \node[thActive,minimum size=0.50cm] at (rtwo) {$2$};
    \node[thInactive,minimum size=0.50cm] at (rthree) {$3$};
    \node[thActive,minimum size=0.50cm] at (rfour) {$4$};
    \node[thActive,minimum size=0.50cm] at (rfive) {$5$};
  \end{scope}
  \node[thLabel,text=ThTeal] at (8.55,-1.30)
    {cost $=\cost(T\vert_A)$};
\end{tikzpicture}}

% --------------------------------------------------------------------------
% A priori master order and realized shortcut.
% --------------------------------------------------------------------------
\newcommand{\ThMasterOrder}{%
\begin{tikzpicture}[x=1cm,y=1cm]
  \node[anchor=west,font=\fontsize{9pt}{11pt}\selectfont\bfseries,text=ThInk]
    at (0,2.65) {Order fixed before activation};
  \node[thDepot] (mr) at (0.7,2.0) {$r$};
  \node[thActive] (mone) at (2.0,2.0) {$1$};
  \node[thInactive] (mtwo) at (3.3,2.0) {$2$};
  \node[thActive] (mthree) at (4.6,2.0) {$3$};
  \node[thInactive] (mfour) at (5.9,2.0) {$4$};
  \node[thInactive] (mfive) at (7.2,2.0) {$5$};
  \node[thActive] (msix) at (8.5,2.0) {$6$};
  \node[thDepot] (mrb) at (9.8,2.0) {$r$};
  \foreach \a/\b in {r/one,one/two,two/three,three/four,four/five,five/six,six/rb}{
    \draw[thFaint,-{Stealth[length=1.4mm]}] (m\a)--(m\b);
  }
  \node[thSmallLabel,anchor=north] at (3.3,1.63) {inactive};
  \node[thSmallLabel,anchor=north] at (5.9,1.63) {inactive};
  \node[thSmallLabel,anchor=north] at (7.2,1.63) {inactive};

  \node[anchor=west,font=\fontsize{9pt}{11pt}\selectfont\bfseries,text=ThInk]
    at (0,0.93) {Realized tour};
  \node[thDepot] (rr) at (0.7,0.28) {$r$};
  \node[thActive] (a1) at (2.8,0.28) {$1$};
  \node[thActive] (a3) at (5.15,0.28) {$3$};
  \node[thActive] (a6) at (7.55,0.28) {$6$};
  \node[thDepot] (rr2) at (9.8,0.28) {$r$};
  \draw[thArrow] (rr)--node[above,thSmallLabel] {$d(r,1)$} (a1);
  \draw[thArrow] (a1)--node[above,thSmallLabel] {$d(1,3)$} (a3);
  \draw[thArrow] (a3)--node[above,thSmallLabel] {$d(3,6)$} (a6);
  \draw[thArrow] (a6)--(rr2);
  \node[thLabel,anchor=west,text=ThTeal] at (10.22,0.28)
    {inactive vertices\\are shortcut};
\end{tikzpicture}}

% --------------------------------------------------------------------------
% Three information regimes.
% --------------------------------------------------------------------------
\newcommand{\ThBenchmarkLadder}{%
\begin{tikzpicture}[x=1cm,y=1cm]
  \draw[ThLine,line width=2.0pt] (0.7,0)--(11.1,0);
  \draw[ThTeal,line width=2.8pt] (0.7,0)--(3.55,0);
  \draw[ThBlue,line width=2.8pt] (3.55,0)--(7.25,0);
  \draw[ThGold,line width=2.8pt] (7.25,0)--(11.1,0);
  \foreach \x/\col/\num in {1.05/ThTeal/1,5.4/ThBlue/2,10.5/ThGold/3}{
    \node[circle,fill=\col,text=ThWhite,minimum size=0.52cm,
      font=\fontsize{10pt}{11pt}\selectfont\bfseries] at (\x,0) {\num};
  }
  \node[anchor=north west,text width=3.25cm,align=left] at (0.55,2.70) {%
    {\fontsize{14pt}{16pt}\selectfont\bfseries A posteriori}\par
    \vspace{0.12cm}
    {\fontsize{10.2pt}{12.2pt}\selectfont\color{ThInkSoft}
      Observe all activations,\\then choose a tour.}};
  \node[anchor=north,text width=3.55cm,align=center] at (5.4,2.70) {%
    {\fontsize{14pt}{16pt}\selectfont\bfseries Adaptive}\par
    \vspace{0.12cm}
    {\fontsize{10.2pt}{12.2pt}\selectfont\color{ThInkSoft}
      Probe sequentially;\\move when active.}};
  \node[anchor=north east,text width=3.35cm,align=right] at (11.0,2.70) {%
    {\fontsize{14pt}{16pt}\selectfont\bfseries A priori}\par
    \vspace{0.12cm}
    {\fontsize{10.2pt}{12.2pt}\selectfont\color{ThInkSoft}
      Fix one order before\\seeing activations.}};
  \node[anchor=north west,text=ThTeal,
    font=\fontsize{9.5pt}{11.5pt}\selectfont\bfseries] at (0.62,-0.38)
    {full information};
  \node[anchor=north,text=ThBlue,
    font=\fontsize{9.5pt}{11.5pt}\selectfont\bfseries] at (5.4,-0.38)
    {information revealed over time};
  \node[anchor=north east,text=ThGold!76!ThInk,
    font=\fontsize{9.5pt}{11.5pt}\selectfont\bfseries] at (11.0,-0.38)
    {no activation information};
\end{tikzpicture}}

% --------------------------------------------------------------------------
% Two realizations of one adaptive probe, shown on one page.
% --------------------------------------------------------------------------
\newcommand{\ThProbeRealizations}{%
\begin{tikzpicture}[x=1cm,y=1cm]
  \path[use as bounding box] (0,-0.50) rectangle (11,1.80);
  \coordinate (lu) at (1.25,0.45);
  \coordinate (lv) at (4.05,0.45);
  \coordinate (ru) at (6.95,0.45);
  \coordinate (rv) at (9.75,0.45);

  \draw[ThLine,line width=0.75pt] (lu)--(lv);
  \draw[ThLine,line width=0.75pt] (ru)--(rv);
  \draw[thArrow,draw=ThTeal,line width=1.55pt]
    (1.62,0.45)--(3.67,0.45);

  \node[circle,draw=ThBlue,fill=ThPaper,line width=1.0pt,
    minimum size=0.72cm,inner sep=0pt,
    font=\fontsize{10.5pt}{12pt}\selectfont] at (lu) {$u$};
  \node[thActive,minimum size=0.72cm,
    font=\fontsize{10.5pt}{12pt}\selectfont] at (lv) {$v$};
  \node[circle,draw=ThBlue,fill=ThPaper,line width=1.0pt,
    minimum size=0.72cm,inner sep=0pt,
    font=\fontsize{10.5pt}{12pt}\selectfont] at (ru) {$u$};
  \node[thInactive,minimum size=0.72cm,
    font=\fontsize{10.5pt}{12pt}\selectfont] at (rv) {$v$};

  \node[thLabel] at (1.25,-0.22) {start at \(u\)};
  \node[thLabel] at (4.05,-0.22) {probe \(v\)};
  \node[thLabel] at (6.95,-0.22) {start at \(u\)};
  \node[thLabel] at (9.75,-0.22) {probe \(v\)};

  % \node[thSmallLabel,text=ThGold!72!ThInk] at (4.05,1.52)
  %   {current position};
  \fill[ThGold!78!ThInk]
    (4.05,0.96)--(3.88,1.27)--(4.22,1.27)--cycle;
  % \node[thSmallLabel,text=ThGold!72!ThInk] at (6.95,1.52)
    % {current position};
  \fill[ThGold!78!ThInk]
    (6.95,0.96)--(6.78,1.27)--(7.12,1.27)--cycle;

\end{tikzpicture}}

% --------------------------------------------------------------------------
% Pairwise gaps among the three information benchmarks.
% --------------------------------------------------------------------------
\newcommand{\ThGapArcs}{%
\begin{tikzpicture}[x=1cm,y=1cm]
  \path[use as bounding box] (0,-2.35) rectangle (11,1.40);
  \coordinate (post) at (1.15,0);
  \coordinate (adapt) at (5.50,0);
  \coordinate (prior) at (9.85,0);

  % Draw the outer arc first so the two local comparisons remain foregrounded.
  \draw[ThGold!80!ThInk,line width=1.25pt,line cap=round]
    (post) to[bend right=30]
    node[midway,below=3pt,fill=ThPaper,inner sep=1.5pt,
      text=ThGold!80!ThInk,
      font=\fontsize{9.4pt}{11.4pt}\selectfont\bfseries]
      {obliviousness gap} (prior);
  \draw[ThTeal,line width=1.25pt,line cap=round]
    (post) to[bend left=23]
    node[midway,above=3pt,fill=ThPaper,inner sep=1.5pt,text=ThTeal,
      font=\fontsize{9.4pt}{11.4pt}\selectfont\bfseries]
      {clairvoyance gap} (adapt);
  \draw[ThBlue,line width=1.25pt,line cap=round]
    (adapt) to[bend left=23]
    node[midway,above=3pt,fill=ThPaper,inner sep=1.5pt,text=ThBlue,
      font=\fontsize{9.4pt}{11.4pt}\selectfont\bfseries]
      {adaptivity gap} (prior);

  \node[fill=ThPaper,inner sep=4pt,
    font=\fontsize{16pt}{19pt}\selectfont] at (post)
    {\(\optpost(\mathcal I)\)};
  \node[fill=ThPaper,inner sep=4pt,
    font=\fontsize{16pt}{19pt}\selectfont] at (adapt)
    {\(\optadapt(\mathcal I)\)};
  \node[fill=ThPaper,inner sep=4pt,
    font=\fontsize{16pt}{19pt}\selectfont] at (prior)
    {\(\optprior(\mathcal I)\)};
  \node[fill=ThPaper,inner sep=2pt,
    font=\fontsize{14pt}{17pt}\selectfont] at (3.325,0)
    {\(\leq\)};
  \node[fill=ThPaper,inner sep=2pt,
    font=\fontsize{14pt}{17pt}\selectfont] at (7.675,0)
    {\(\leq\)};
\end{tikzpicture}}

% --------------------------------------------------------------------------
% Three-chain construction.
% --------------------------------------------------------------------------
\newcommand{\ThThreeChain}{%
\begin{tikzpicture}[x=0.93cm,y=0.93cm]
  \node[thDepot] (A0) at (0, 2.35) {$a_0$};
  \node[thDet] (A1) at (1.4, 2.35) {$a_1$};
  \node[thDet] (Am) at (6.0, 2.35) {$a_m$};
  \node[thDet] (Bm) at (0, 0) {$b_m$};
  \node[thDet] (B1) at (4.6, 0) {$b_1$};
  \node[thDet] (B0) at (6.0, 0) {$b_0$};
  \node[thDet] (C0) at (0,-2.35) {$c_0$};
  \node[thDet] (C1) at (1.4,-2.35) {$c_1$};
  \node[thDet] (Cm) at (6.0,-2.35) {$c_m$};
  \node[thStoch,minimum size=0.60cm] (S) at (7.35,-1.18) {$s$};

  \draw[ThMuted,densely dashed,line width=0.8pt] (A0)--(Bm);
  \draw[ThMuted,densely dashed,line width=0.8pt] (A0)--(Cm);
  \draw[ThMuted,densely dashed,line width=0.8pt] (Am)--(B0);
  \draw[ThMuted,densely dashed,line width=0.8pt] (Am)--(C0);
  \draw[ThCoral,dash pattern=on 4pt off 2pt on 1pt off 2pt,line width=0.8pt]
    (Bm)--(C0);
  \draw[ThCoral,dash pattern=on 4pt off 2pt on 1pt off 2pt,line width=0.8pt]
    (B0)--(S);
  \draw[ThCoral,dash pattern=on 4pt off 2pt on 1pt off 2pt,line width=0.8pt]
    (Cm)--(S);
  \draw[ThBlue,line width=1.15pt] (A0)--(A1)--(Am);
  \draw[ThBlue,line width=1.15pt] (Bm)--(B1)--(B0);
  \draw[ThBlue,line width=1.15pt] (C0)--(C1)--(Cm);
  \node[fill=ThPaper,inner sep=3pt,text=ThMuted] at (3.55,2.35) {$\cdots$};
  \node[fill=ThPaper,inner sep=3pt,text=ThMuted] at (2.55,0) {$\cdots$};
  \node[fill=ThPaper,inner sep=3pt,text=ThMuted] at (3.55,-2.35) {$\cdots$};
  \node[thLabel,anchor=east,font=\fontsize{10pt}{12pt}\selectfont\bfseries]
    at (-0.55,2.35) {$\mathbf{A}$};
  \node[thLabel,anchor=east,font=\fontsize{10pt}{12pt}\selectfont\bfseries]
    at (-0.55,0) {$\mathbf{B}$};
  \node[thLabel,anchor=east,font=\fontsize{10pt}{12pt}\selectfont\bfseries]
    at (-0.55,-2.35) {$\mathbf{C}$};
  \node[thSmallLabel,anchor=west,text=ThBlue] at (7.0,1.55)
    {chain edges: \(2\)};
  \node[thSmallLabel,anchor=west,text=ThMuted] at (7.0,1.05)
    {dashed: \(U\)};
  \node[thSmallLabel,anchor=west,text=ThCoral] at (7.0,0.55)
    {dash-dotted: \(P\)};
\end{tikzpicture}}

\newcommand{\ThOrderConflict}{%
\begin{tikzpicture}[x=1cm,y=1cm]
  \path[use as bounding box] (0,-0.45) rectangle (11,1.55);
  \node[thDet] (A0) at (0.55,0.85) {$A$};
  \node[thDet] (B0) at (2.25,0.85) {$B$};
  \node[thDet] (C0) at (3.95,0.85) {$C$};
  \draw[thArrow,draw=ThTeal] (A0)--(B0);
  \draw[thArrow,draw=ThTeal] (B0)--(C0);
  \draw[thArrow,draw=ThTeal] (C0) to[bend left=33] (A0);
  \draw[ThLine,line width=0.7pt] (5.15,-0.45)--(5.15,1.55);
  \node[thDet] (A1) at (6.05,0.85) {$A$};
  \node[thDet] (C1) at (7.55,0.85) {$C$};
  \node[thActive] (S1) at (9.05,0.85) {$s$};
  \node[thDet] (B1) at (10.55,0.85) {$B$};
  \draw[thArrow,draw=ThCoral] (A1)--(C1);
  \draw[thArrow,draw=ThCoral] (C1)--(S1);
  \draw[thArrow,draw=ThCoral] (S1)--(B1);
  \draw[thArrow,draw=ThCoral] (B1) to[bend left=22] (A1);
\end{tikzpicture}}

% --------------------------------------------------------------------------
% Five-vertex symmetric adaptivity-gap demo.
% --------------------------------------------------------------------------
\newcommand{\ThKFourSubdivision}{%
\begin{tikzpicture}[x=1cm,y=1cm]
  \path[use as bounding box] (-2.70,-1.90) rectangle (2.70,2.35);
  \coordinate (r) at (-2.05,0.55);
  \coordinate (a) at ( 2.05,0.55);
  \coordinate (b) at (-2.05,-1.25);
  \coordinate (c) at ( 2.05,-1.25);
  \coordinate (s) at ( 0,1.72);

  % Unit-edge K_4 with ra replaced by the two-edge path r-s-a.
  \foreach \u/\v in {r/b,a/c,b/c,r/c,a/b}{
    \draw[thEdge] (\u)--(\v);
  }
  \draw[ThBlue,line width=1.55pt,line join=round] (r)--(s)--(a);

  \node[thDepot,minimum size=0.62cm,
    font=\fontsize{10pt}{11pt}\selectfont] at (r) {$r$};
  \node[thDet,minimum size=0.62cm,
    font=\fontsize{10pt}{11pt}\selectfont] at (a) {$a$};
  \node[thDet,minimum size=0.62cm,
    font=\fontsize{10pt}{11pt}\selectfont] at (b) {$b$};
  \node[thDet,minimum size=0.62cm,
    font=\fontsize{10pt}{11pt}\selectfont] at (c) {$c$};
  \node[thStoch,minimum size=0.64cm,
    font=\fontsize{10pt}{11pt}\selectfont] at (s) {$s$};
  \node[thLabel,text=ThTeal,anchor=west] at (0.35,1.72)
    {$p_s=\tfrac12$};
\end{tikzpicture}}

\newcommand{\ThKFourOutcomes}{%
\begin{tikzpicture}[x=1cm,y=1cm]
  \path[use as bounding box] (0,-2.10) rectangle (11,1.90);
  \coordinate (r0) at (0.70,0.45);
  \coordinate (a0) at (4.40,0.45);
  \coordinate (b0) at (0.70,-0.95);
  \coordinate (c0) at (4.40,-0.95);
  \coordinate (s0) at (2.55,1.48);
  \foreach \u/\v in {r0/b0,a0/c0,b0/c0,r0/c0,a0/b0,r0/s0,s0/a0}{
    \draw[thFaint] (\u)--(\v);
  }
  \draw[thArrowThin,draw=ThBlue] (r0)--(b0);
  \draw[thArrowThin,draw=ThBlue] (b0)--(a0);
  \draw[thArrowThin,draw=ThBlue] (a0)--(c0);
  \draw[thArrowThin,draw=ThBlue] (c0)--(r0);
  \node[thDepot] at (r0) {$r$};
  \node[thDet] at (a0) {$a$};
  \node[thDet] at (b0) {$b$};
  \node[thDet] at (c0) {$c$};
  \node[thInactive] at (s0) {$s$};
  \node[font=\fontsize{10pt}{12pt}\selectfont,text=ThBlue]
    at (2.55,-1.65) {$r\to b\to a\to c\to r$: cost \(4\)};

  \draw[ThLine,line width=0.7pt] (5.50,-1.85)--(5.50,1.85);
  \coordinate (r1) at (6.60,0.45);
  \coordinate (a1) at (10.30,0.45);
  \coordinate (b1) at (6.60,-0.95);
  \coordinate (c1) at (10.30,-0.95);
  \coordinate (s1) at (8.45,1.48);
  \foreach \u/\v in {r1/b1,a1/c1,b1/c1,r1/c1,a1/b1,r1/s1,s1/a1}{
    \draw[thFaint] (\u)--(\v);
  }
  \draw[thArrowThin,draw=ThTeal] (r1)--(s1);
  \draw[thArrowThin,draw=ThTeal] (s1)--(a1);
  \draw[thArrowThin,draw=ThTeal] (a1)--(c1);
  \draw[thArrowThin,draw=ThTeal] (c1)--(b1);
  \draw[thArrowThin,draw=ThTeal] (b1)--(r1);
  \node[thDepot] at (r1) {$r$};
  \node[thDet] at (a1) {$a$};
  \node[thDet] at (b1) {$b$};
  \node[thDet] at (c1) {$c$};
  \node[thActive] at (s1) {$s$};
  \node[font=\fontsize{10pt}{12pt}\selectfont,text=ThTeal]
    at (8.45,-1.65) {$r\to s\to a\to c\to b\to r$: cost \(5\)};
\end{tikzpicture}}

% --------------------------------------------------------------------------
% Toy probabilistic-method experiment: two equally likely (1,2)-metrics.
% --------------------------------------------------------------------------
\newcommand{\ThRandomMetricToy}{%
\begin{tikzpicture}[x=1cm,y=1cm]
  \path[use as bounding box] (0,-1.10) rectangle (11,1.25);
  \draw[ThLine,line width=0.65pt] (5.50,-1.05)--(5.50,1.20);
  \coordinate (lr) at (1.20,0.90);
  \coordinate (la) at (4.20,0.90);
  \coordinate (lb) at (4.20,-0.75);
  \coordinate (lc) at (1.20,-0.75);
  \draw[ThGold!78!ThInk,densely dashed,line width=0.85pt] (lr)--(lb);
  \draw[ThGold!78!ThInk,densely dashed,line width=0.85pt] (la)--(lc);
  \draw[ThTeal,line width=1.75pt,line join=round]
    (lr)--(la)--(lb)--(lc)--cycle;
  \node[thDepot,minimum size=0.52cm] at (lr) {$r$};
  \node[thDet,minimum size=0.48cm] at (la) {$a$};
  \node[thDet,minimum size=0.48cm] at (lb) {$b$};
  \node[thDet,minimum size=0.48cm] at (lc) {$c$};
  \coordinate (rr) at (6.80,0.90);
  \coordinate (ra) at (9.80,0.90);
  \coordinate (rb) at (9.80,-0.75);
  \coordinate (rc) at (6.80,-0.75);
  \draw[ThGold!78!ThInk,densely dashed,line width=0.85pt] (rr)--(rc);
  \draw[ThGold!78!ThInk,densely dashed,line width=0.85pt] (ra)--(rb);
  \draw[ThTeal,line width=1.75pt,line join=round]
    (rr)--(ra)--(rc)--(rb)--cycle;
  \node[thDepot,minimum size=0.52cm] at (rr) {$r$};
  \node[thDet,minimum size=0.48cm] at (ra) {$a$};
  \node[thDet,minimum size=0.48cm] at (rb) {$b$};
  \node[thDet,minimum size=0.48cm] at (rc) {$c$};
\end{tikzpicture}}

% --------------------------------------------------------------------------
% The same tour is evaluated under both toy metrics.
% --------------------------------------------------------------------------
\newcommand{\ThRandomMetricToyComparison}{%
\begin{tikzpicture}[x=1cm,y=1cm]
  \path[use as bounding box] (0,-1.10) rectangle (11,1.85);
  \node[font=\fontsize{10.5pt}{12.5pt}\selectfont\bfseries,text=ThInkSoft]
    at (5.20,1.60) {fix \(\boldsymbol{\mathbf{T=(r,a,b,c,r)}}\) before seeing \(\mathbf{D}\)};

  \foreach \x/\name in {1.45/lr,2.75/la,4.05/lb,5.35/lc,6.65/lrr}{
    \coordinate (\name) at (\x,0.72);
  }
  \foreach \u/\v in {lr/la,la/lb,lb/lc,lc/lrr}{
    \draw[ThTeal,line width=1.45pt] (\u)--(\v);
  }
  \node[thDepot,minimum size=0.43cm] at (lr) {$r$};
  \node[thDet,minimum size=0.40cm] at (la) {$a$};
  \node[thDet,minimum size=0.40cm] at (lb) {$b$};
  \node[thDet,minimum size=0.40cm] at (lc) {$c$};
  \node[thDepot,minimum size=0.43cm] at (lrr) {$r$};
  \node[anchor=east,font=\fontsize{10pt}{12pt}\selectfont\bfseries,
    text=ThInkSoft] at (1.05,0.72) {under \(\mathbf{D_1}\)};
  \node[anchor=west,font=\fontsize{10pt}{12pt}\selectfont,text=ThTeal!72!ThInk]
    at (7.05,0.72) {cost \(1+1+1+1=4\)};

  \foreach \x/\name in {1.45/br,2.75/ba,4.05/bb,5.35/bc,6.65/brr}{
    \coordinate (\name) at (\x,-0.38);
  }
  \draw[ThTeal,line width=1.45pt] (br)--(ba);
  \draw[ThGold,line width=1.45pt] (ba)--(bb);
  \draw[ThTeal,line width=1.45pt] (bb)--(bc);
  \draw[ThGold,line width=1.45pt] (bc)--(brr);
  \node[thDepot,minimum size=0.43cm] at (br) {$r$};
  \node[thDet,minimum size=0.40cm] at (ba) {$a$};
  \node[thDet,minimum size=0.40cm] at (bb) {$b$};
  \node[thDet,minimum size=0.40cm] at (bc) {$c$};
  \node[thDepot,minimum size=0.43cm] at (brr) {$r$};
  \node[anchor=east,font=\fontsize{10pt}{12pt}\selectfont\bfseries,
    text=ThInkSoft] at (1.05,-0.38) {under \(\mathbf{D_2}\)};
  \node[anchor=west,font=\fontsize{10pt}{12pt}\selectfont,text=ThGold!82!ThInk]
    at (7.05,-0.38) {cost \(1+2+1+2=6\)};

  \node[font=\fontsize{10.5pt}{12.5pt}\selectfont\bfseries,text=ThBlue]
    at (9.45,0.17) {expected cost \(\boldsymbol{\mathbf{=5}}\)};
\end{tikzpicture}}

% --------------------------------------------------------------------------
% One realized a-priori shortcut pair in the random (1,2)-metric.
% --------------------------------------------------------------------------
\newcommand{\ThRandomShortcutPair}{%
\begin{tikzpicture}[x=1cm,y=1cm]
  \path[use as bounding box] (-0.20,0.50) rectangle (7.65,2.05);

  \coordinate (u0) at (1.05,0.92);
  \coordinate (v10) at (2.40,0.92);
  \coordinate (v20) at (3.75,0.92);
  \coordinate (dots0) at (5.02,0.92);
  \coordinate (w0) at (6.32,0.92);
  \coordinate (dots1) at (4.42,0.92);
  \coordinate (dots2) at (5.62,0.92);

  \draw[ThMuted,line width=0.72pt] (0.42,0.92)--(u0);
  \draw[thArrowThin,draw=ThMuted]
    (u0)--(v10)--(v20);
  \draw[thArrowThin,draw=ThMuted]
    (v20)--(dots1);
  \draw[thArrowThin,draw=ThMuted]
    (dots2)--(w0)--(7.00,0.92);
  \node[font=\fontsize{13pt}{14pt}\selectfont,text=ThMuted]
    (leftdots) at (0.10,0.92) {$\cdots$};
  \node[thStoch,minimum size=0.45cm] (unode) at (u0) {$u$};
  \node[thStoch,minimum size=0.45cm] at (v10) {$v_1$};
  \node[thStoch,minimum size=0.45cm] at (v20) {$v_2$};
  \node[font=\fontsize{13pt}{14pt}\selectfont,text=ThMuted] at (dots0) {$\cdots$};
  \node[thStoch,minimum size=0.45cm] (wnode) at (w0) {$w$};
  \node[font=\fontsize{13pt}{14pt}\selectfont,text=ThMuted]
    (rightdots) at (7.40,0.92) {$\cdots$};

  \node[fit=(leftdots)(rightdots)(unode)(wnode),inner sep=0pt]
    (tourExcerpt) {};
  % \node[thLabel,font=\fontsize{10pt}{12pt}\selectfont\bfseries,
  %   text=ThInkSoft,above=4mm of tourExcerpt]
  %   {For a fixed a priori tour \(T\)};
\end{tikzpicture}}

% --------------------------------------------------------------------------
% Random (1,2)-metric intuition.
% --------------------------------------------------------------------------
\newcommand{\ThRandomMetric}{%
\begin{tikzpicture}[x=1cm,y=1cm]
  \path[use as bounding box] (-2.65,-2.35) rectangle (2.65,2.35);
  \foreach \i in {1,...,10}{
    \coordinate (v\i) at ({36*\i}:2.05);
  }

  % Draw the sampled short-edge graph. Omitted pairs have metric length 2.
  \foreach \a/\b in {1/2,1/4,1/7,2/5,2/9,3/5,3/8,4/6,4/10,
                       5/7,5/10,6/8,7/9,8/10}{
    \draw[ThTeal!78!ThInk,line width=0.82pt] (v\a)--(v\b);
  }
  \foreach \i in {1,2,3,4,6,7,8,9,10}{
    \node[thStoch,minimum size=0.42cm,inner sep=0.5pt] at (v\i) {};
  }
  \node[thDepot,minimum size=0.46cm] at (v5) {$r$};
\end{tikzpicture}}

% --------------------------------------------------------------------------
% The active vertices inherit an Erdos--Renyi short-edge graph.
% --------------------------------------------------------------------------
\newcommand{\ThActiveInducedHamiltonian}{%
\begin{tikzpicture}[x=1cm,y=1cm]
  \path[use as bounding box] (-2.65,-2.35) rectangle (2.65,2.35);
  \coordinate (r) at ({180}:2.02);
  \foreach \i/\angle in {1/128,2/76,3/24,4/-28,5/-80,6/-132}{
    \coordinate (a\i) at ({\angle}:2.02);
  }

  % Additional short edges are shown faintly; one all-short tour is emphasized.
  \foreach \u/\v in {r/a2,r/a4,a1/a4,a1/a5,a2/a5,a2/a6,a3/a6,a3/a5}{
    \draw[ThLine,line width=0.62pt] (\u)--(\v);
  }
  \draw[ThTeal,line width=1.75pt,line join=round]
    (r)--(a1)--(a2)--(a3)--(a4)--(a5)--(a6)--cycle;

  \node[thDepot,minimum size=0.50cm] at (r) {$r$};
  \foreach \i in {1,...,6}{
    \node[thActive,minimum size=0.46cm] at (a\i) {$\i$};
  }
\end{tikzpicture}}

% --------------------------------------------------------------------------
% A fixed order becomes mostly long after inactive vertices are skipped.
% --------------------------------------------------------------------------
\newcommand{\ThFixedOrderShortcutWeight}{%
\begin{tikzpicture}[x=1cm,y=1cm]
  \path[use as bounding box] (0,-1.05) rectangle (11,2.10);

  \foreach \x/\name in {0.65/tzero,1.85/tone,3.05/ttwo,4.25/tthree,
                          5.45/tfour,6.65/tfive,7.85/tsix,9.05/tseven,
                          10.25/teight}{
    \coordinate (\name) at (\x,1.15);
  }
  \draw[thArrowThin,draw=ThMuted]
    (tzero)--(tone)--(ttwo)--(tthree)--(tfour)--(tfive)--(tsix)--(tseven)--(teight);
  \node[thDepot,minimum size=0.48cm] at (tzero) {$r$};
  \node[thDepot,minimum size=0.48cm] at (teight) {$r$};
  \foreach \name in {tone,ttwo,tthree,tfour,tfive,tsix,tseven}{
    \node[thStoch,minimum size=0.45cm] at (\name) {};
  }
  \node[anchor=east,font=\fontsize{10pt}{12pt}\selectfont\bfseries,
    text=ThInkSoft] at (0.25,1.15) {fixed order};

  \coordinate (bzero) at (0.65,-0.25);
  \coordinate (btwo) at (3.05,-0.25);
  \coordinate (bfive) at (6.65,-0.25);
  \coordinate (bseven) at (9.05,-0.25);
  \coordinate (beight) at (10.25,-0.25);
  \draw[ThGold,line width=1.45pt]
    (bzero)--node[midway,above=2pt,font=\fontsize{9pt}{10pt}\selectfont,
      text=ThGold!80!ThInk] {2} (btwo);
  \draw[ThGold,line width=1.45pt]
    (btwo)--node[midway,above=2pt,font=\fontsize{9pt}{10pt}\selectfont,
      text=ThGold!80!ThInk] {2} (bfive);
  \draw[ThTeal,line width=1.45pt]
    (bfive)--node[midway,above=2pt,font=\fontsize{9pt}{10pt}\selectfont,
      text=ThTeal!72!ThInk] {1} (bseven);
  \draw[ThGold,line width=1.45pt]
    (bseven)--node[midway,above=2pt,font=\fontsize{9pt}{10pt}\selectfont,
      text=ThGold!80!ThInk] {2} (beight);
  \node[thDepot,minimum size=0.48cm] at (bzero) {$r$};
  \node[thActive,minimum size=0.45cm] at (btwo) {};
  \node[thActive,minimum size=0.45cm] at (bfive) {};
  \node[thActive,minimum size=0.45cm] at (bseven) {};
  \node[thDepot,minimum size=0.48cm] at (beight) {$r$};
  \node[anchor=east,font=\fontsize{10pt}{12pt}\selectfont\bfseries,
    text=ThInkSoft] at (0.25,-0.25) {one shortcut};
  \node[thSmallLabel,text=ThGold!82!ThInk] at (5.50,-0.82)
    {most realized jumps have length \(2\)};

\end{tikzpicture}}

% --------------------------------------------------------------------------
% Local view of the nearest-unprobed-neighbor adaptive policy.
% --------------------------------------------------------------------------
\newcommand{\ThShortNeighborPolicy}{%
\begin{tikzpicture}[x=1cm,y=1cm]
  \path[use as bounding box] (-0.40,-1.70) rectangle (7.00,1.70);
  \coordinate (u) at (0.65,0.00);
  \coordinate (sone) at (3.40,1.35);
  \coordinate (stwo) at (3.40,0.00);
  \coordinate (sthree) at (3.40,-1.35);
  \coordinate (fallback) at (6.25,0.00);

  % Draw available length-1 edges and the rare length-2 fallback first.
  \foreach \v in {sone,stwo,sthree}{
    \draw[ThTeal!78!ThInk,line width=1.05pt] (u)--(\v);
  }
  \draw[ThMuted,densely dashed,line width=0.85pt] (u)--(fallback);
  \draw[thArrow,draw=ThTeal] (u)--(sthree);

  \node[thActive,minimum size=0.62cm] at (u) {$u$};
  \node[thInactive,minimum size=0.52cm] at (sone) {};
  \node[thInactive,minimum size=0.52cm] at (stwo) {};
  \node[thActive,minimum size=0.52cm] at (sthree) {$v$};
  \node[thStoch,minimum size=0.52cm] at (fallback) {};

  \node[thSmallLabel] at (6.25,-0.70) {length-\(2\) fallback};
\end{tikzpicture}}

% --------------------------------------------------------------------------
% Same number of moves, but two different typical movement lengths.
% --------------------------------------------------------------------------
\newcommand{\ThFactorTwoTourComparison}{%
\begin{tikzpicture}[x=1cm,y=1cm]
  \path[use as bounding box] (0,-1.15) rectangle (11,1.75);
  \foreach \x/\i in {1.45/0,2.75/1,4.05/2,5.35/3,6.65/4,7.95/5}{
    \coordinate (p\i) at (\x,0.95);
    \coordinate (a\i) at (\x,-0.35);
  }
  \foreach \i/\j in {0/1,1/2,2/3,3/4,4/5}{
    \draw[ThGold,line width=1.55pt] (p\i)--(p\j)
      node[midway,above=1.5pt,font=\fontsize{9pt}{10pt}\selectfont,
        text=ThGold!82!ThInk] {2};
    \draw[ThTeal,line width=1.55pt] (a\i)--(a\j)
      node[midway,above=1.5pt,font=\fontsize{9pt}{10pt}\selectfont,
        text=ThTeal!72!ThInk] {1};
  }
  \foreach \i in {0,...,5}{
    \node[circle,draw=ThInkSoft,fill=ThWhite,line width=0.75pt,
      minimum size=0.30cm,inner sep=0pt] at (p\i) {};
    \node[circle,draw=ThTeal!55!ThInk,fill=ThTeal,line width=0.75pt,
      minimum size=0.30cm,inner sep=0pt] at (a\i) {};
  }
  \node[anchor=east,font=\fontsize{10pt}{12pt}\selectfont\bfseries,
    text=ThInkSoft] at (1.02,0.95) {a priori};
  \node[anchor=east,font=\fontsize{10pt}{12pt}\selectfont\bfseries,
    text=ThInkSoft] at (1.02,-0.35) {adaptive};
  \node[anchor=west,font=\fontsize{11pt}{13pt}\selectfont\bfseries,
    text=ThGold!82!ThInk] at (8.45,0.95) {almost \(\mathbf{2}\) per move};
  \node[anchor=west,font=\fontsize{11pt}{13pt}\selectfont\bfseries,
    text=ThTeal!72!ThInk] at (8.45,-0.35) {almost \(\mathbf{1}\) per move};
\end{tikzpicture}}

% --------------------------------------------------------------------------
% Why the same random metric separates a priori and adaptive behavior.
% --------------------------------------------------------------------------
\newcommand{\ThFactorTwoMechanisms}{%
\begin{tikzpicture}[x=1cm,y=1cm]
  \path[use as bounding box] (0,-1.25) rectangle (11,2.10);
  \draw[ThLine,line width=0.65pt] (5.50,-1.05)--(5.50,1.95);

  % A priori: activation turns a carefully chosen order into new shortcut pairs.
  \node[font=\fontsize{10.5pt}{12.5pt}\selectfont\bfseries,text=ThInkSoft]
    at (2.75,1.82) {A priori: fix one order};
  \foreach \x/\name in {0.55/lr,1.42/lone,2.29/ltwo,3.16/lthree,
                           4.03/lfour,4.90/lrr}{
    \coordinate (\name) at (\x,0.92);
  }
  \draw[thArrowThin,draw=ThMuted]
    (lr)--(lone)--(ltwo)--(lthree)--(lfour)--(lrr);
  \node[thDepot,minimum size=0.40cm] at (lr) {$r$};
  \node[thDepot,minimum size=0.40cm] at (lrr) {$r$};
  \foreach \name in {lone,ltwo,lthree,lfour}{
    \node[thStoch,minimum size=0.38cm] at (\name) {};
  }

  \foreach \x/\name in {0.55/br,1.70/bone,3.10/btwo,4.35/bthree,4.90/brr}{
    \coordinate (\name) at (\x,-0.12);
  }
  \draw[ThGold,line width=1.30pt] (br)--
    node[midway,above=1.5pt,font=\fontsize{8pt}{9pt}\selectfont,
      text=ThGold!82!ThInk] {2} (bone);
  \draw[ThGold,line width=1.30pt] (bone)--
    node[midway,above=1.5pt,font=\fontsize{8pt}{9pt}\selectfont,
      text=ThGold!82!ThInk] {2} (btwo);
  \draw[ThTeal,line width=1.30pt] (btwo)--
    node[midway,above=1.5pt,font=\fontsize{8pt}{9pt}\selectfont,
      text=ThTeal!72!ThInk] {1} (bthree);
  \draw[ThGold,line width=1.30pt] (bthree)--
    node[midway,above=1.5pt,font=\fontsize{8pt}{9pt}\selectfont,
      text=ThGold!82!ThInk] {2} (brr);
  \node[thDepot,minimum size=0.40cm] at (br) {$r$};
  \node[thDepot,minimum size=0.40cm] at (brr) {$r$};
  \foreach \name in {bone,btwo,bthree}{
    \node[thActive,minimum size=0.38cm] at (\name) {};
  }
  \node[thSmallLabel,text=ThGold!82!ThInk] at (2.75,-0.78)
    {activation creates new, usually long shortcut pairs};

  % Adaptive: the current location can test its known short neighbors first.
  \node[font=\fontsize{10.5pt}{12.5pt}\selectfont\bfseries,text=ThInkSoft]
    at (8.25,1.82) {Adaptive: react to outcomes};
  \coordinate (u) at (6.45,0.20);
  \coordinate (inactive) at (8.25,1.00);
  \coordinate (active) at (8.25,-0.58);
  \coordinate (fallback) at (10.20,0.20);
  \draw[ThTeal!78!ThInk,line width=1.00pt] (u)--(inactive);
  \draw[ThTeal!78!ThInk,line width=1.00pt] (u)--(active);
  \draw[ThMuted,densely dashed,line width=0.82pt] (u)--(fallback);
  \draw[thArrow,draw=ThTeal] (u)--(active);
  \node[thActive,minimum size=0.50cm] at (u) {$u$};
  \node[thInactive,minimum size=0.46cm] at (inactive) {};
  \node[thActive,minimum size=0.46cm] at (active) {$v$};
  \node[thStoch,minimum size=0.46cm] at (fallback) {};
  \node[thSmallLabel,anchor=west] at (8.63,1.00) {inactive: stay};
  \node[thSmallLabel,anchor=west,text=ThTeal] at (8.63,-0.58)
    {active: visit immediately};
  \node[font=\fontsize{8pt}{9pt}\selectfont,text=ThTeal!72!ThInk]
    at (7.25,0.83) {1};
  \node[font=\fontsize{8pt}{9pt}\selectfont,text=ThTeal!72!ThInk]
    at (7.25,-0.38) {1};
  \node[font=\fontsize{8pt}{9pt}\selectfont,text=ThMuted]
    at (9.22,0.38) {2};
  \node[thSmallLabel] at (10.20,-0.32) {fallback};
  \node[thSmallLabel,text=ThTeal!72!ThInk] at (8.25,-1.02)
    {probe short neighbors first};
\end{tikzpicture}}

% --------------------------------------------------------------------------
% The two high-probability properties overlap for at least one metric.
% --------------------------------------------------------------------------
\newcommand{\ThGoodMetricIntersection}{%
\begin{tikzpicture}[x=1cm,y=1cm]
  \path[use as bounding box] (0,-1.45) rectangle (11,1.65);

  \fill[ThGold!24,fill opacity=0.55]
    (3.85,0.30) ellipse[x radius=2.95,y radius=1.20];
  \draw[ThGold!78!ThInk,line width=1.00pt]
    (3.85,0.30) ellipse[x radius=2.95,y radius=1.20];
  \fill[ThTeal!24,fill opacity=0.55]
    (7.15,0.30) ellipse[x radius=2.95,y radius=1.20];
  \draw[ThTeal!70!ThInk,line width=1.00pt]
    (7.15,0.30) ellipse[x radius=2.95,y radius=1.20];

  \node[align=center,
    font=\fontsize{9.2pt}{11pt}\selectfont,text=ThInk]
    at (2.55,0.38)
    {\textbf{\(\boldsymbol{\mathcal E}_{\mathbf{prior}}\)}\\[2pt]
     \(\optprior(D)\)\\[-1pt]
     \(\ge(2-o(1))np\)};
  \node[align=center,
    font=\fontsize{9.2pt}{11pt}\selectfont,text=ThInk]
    at (8.45,0.38)
    {\textbf{\(\boldsymbol{\mathcal E}_{\mathbf{adapt}}\)}\\[2pt]
     \(\optadapt(D)\)\\[-1pt]
     \(\le(1+o(1))np\)};

  \node[rectangle,rounded corners=1.2pt,draw=ThBlue,fill=ThBlue!8,
    line width=0.9pt,minimum width=1.02cm,minimum height=0.62cm,
    font=\fontsize{9.5pt}{11pt}\selectfont\bfseries,text=ThBlue]
    at (5.50,0.30) {fix \(\mathbf{d_n}\)};

  \node[font=\fontsize{8.7pt}{10pt}\selectfont,text=ThGold!82!ThInk]
    at (2.55,-1.10) {$\Pp_D(\mathcal E_{\mathrm{prior}})=1-o(1)$};
  \node[font=\fontsize{8.7pt}{10pt}\selectfont,text=ThTeal!72!ThInk]
    at (8.45,-1.10) {$\Pp_D(\mathcal E_{\mathrm{adapt}})=1-o(1)$};
\end{tikzpicture}}

% --------------------------------------------------------------------------
% Sampled master tour for the symmetric factor-three upper bound.
% --------------------------------------------------------------------------
\newcommand{\ThSampledBackbone}{%
\begin{tikzpicture}[x=1cm,y=1cm]
  \path[use as bounding box] (-2.65,-2.05) rectangle (2.75,1.75);
  \coordinate (r) at (-2.05,0.10);
  \coordinate (a) at (-0.95,1.45);
  \coordinate (b) at (0.85,1.35);
  \coordinate (c) at (1.95,-0.15);
  \coordinate (d) at (-0.20,-1.45);
  \coordinate (u) at (-2.20,-1.30);
  \coordinate (v) at (2.35,1.25);

  \draw[ThBlue,line width=1.65pt,line join=round]
    (r)--(a)--(b)--(c)--(d)--cycle;
  \draw[ThTeal,line width=0.95pt] ([yshift=0.05cm]u.north east)--([yshift=0.05cm]d.south west);
  \draw[ThTeal,line width=0.95pt] ([yshift=-0.05cm]u.north east)--([yshift=-0.05cm]d.south west);
  \draw[ThTeal,line width=0.95pt] ([yshift=0.05cm]v.south west)--([yshift=0.05cm]b.north east);
  \draw[ThTeal,line width=0.95pt] ([yshift=-0.05cm]v.south west)--([yshift=-0.05cm]b.north east);

  \node[thDepot,minimum size=0.56cm] at (r) {$r$};
  \foreach \x/\lab in {a/1,b/2,c/3,d/4}{
    \node[thDet,minimum size=0.52cm] at (\x) {$\lab$};
  }
  \node[thStoch,minimum size=0.52cm] at (u) {$u$};
  \node[thStoch,minimum size=0.52cm] at (v) {$v$};
\end{tikzpicture}}

% --------------------------------------------------------------------------
% Rectangular cyclic grid and sweep.
% --------------------------------------------------------------------------
\newcommand{\ThRectGrid}{%
\begin{tikzpicture}[>=Stealth,x=1cm,y=1cm]
  \foreach \x/\name/\lab in {0.95/Cone/C_1,3.80/Ctwo/C_2,7.75/Cm/C_{k^2}}{
    \node[draw=ThInkSoft,rounded corners=1pt,minimum width=1.24cm,
      minimum height=2.55cm,fill=ThMist] (\name) at (\x,0) {};
    \node[above=1.5mm of \name,font=\fontsize{10pt}{12pt}\selectfont\bfseries]
      {$\mathbf{\lab}$};
  }

  \node[thDepot,minimum size=0.56cm] at ($(Cone.center)+(0,0.79)$) {$r$};
  \node[thStoch,minimum size=0.36cm] at ($(Cone.center)+(0,0.18)$) {};
  \node[text=ThMuted] at ($(Cone.center)+(0,-0.27)$) {$\vdots$};
  \node[thStoch,minimum size=0.36cm] at ($(Cone.center)+(0,-0.79)$) {};

  \foreach \name in {Ctwo,Cm}{
    \node[thStoch,minimum size=0.39cm] at ($(\name.center)+(0,0.72)$) {};
    \node[text=ThMuted] at (\name.center) {$\vdots$};
    \node[thStoch,minimum size=0.39cm] at ($(\name.center)+(0,-0.72)$) {};
  }
  \node[text=ThMuted,font=\fontsize{15pt}{17pt}\selectfont] at (5.78,0) {$\cdots$};
  \draw[thArrow] (Cone.east)--node[above,thSmallLabel] {unit arcs} (Ctwo.west);
  \draw[thArrow] (Ctwo.east)--(5.28,0);
  \draw[thArrow] (6.28,0)--(Cm.west);
  \draw[thArrow] (Cm.south) to[bend left=32]
    node[midway,above=1.0mm,thSmallLabel] {wraparound} (Cone.south);
\end{tikzpicture}}

\newcommand{\ThColumnSweep}{%
\begin{tikzpicture}[x=1cm,y=1cm]
  \foreach \x/\lab in {0.7/C_1,2.35/C_2,4.0/C_3,6.9/C_j,9.95/C_{k^2}}{
    \node[font=\fontsize{9pt}{11pt}\selectfont\bfseries,text=ThMuted]
      at (\x,2.28) {$\mathbf{\lab}$};
    \draw[ThLine,line width=0.65pt] (\x,-1.25)--(\x,1.88);
  }
  \node[text=ThMuted] at (5.45,2.28) {$\cdots$};
  \node[text=ThMuted] at (8.45,2.28) {$\cdots$};
  \foreach \y/\num/\col in {1.35/1/ThBlue,0.30/2/ThTeal}{
    \node[anchor=east,font=\fontsize{9pt}{11pt}\selectfont\bfseries,text=\col]
      (sweep\num) at (-0.05,\y) {sweep \num};
    \draw[-{Stealth[length=2mm]},draw=\col,line width=1.4pt]
      (0.35,\y)--(10.4,\y);
  }
  \node[below=4mm of sweep2,font=\fontsize{11pt}{13pt}\selectfont,text=ThMuted]
    {$\vdots$};
  \foreach \x/\y/\col in {0.7/1.35/ThBlue,2.35/1.35/ThBlue,4.0/1.35/ThBlue,6.9/1.35/ThBlue,9.95/1.35/ThBlue,
    2.35/0.30/ThTeal,4.0/0.30/ThTeal,6.9/0.30/ThTeal,9.95/0.30/ThTeal}{
    \node[circle,fill=\col,draw=ThWhite,line width=0.5pt,minimum size=0.23cm]
      at (\x,\y) {};
  }
\end{tikzpicture}}

\newcommand{\ThPacketColumns}{%
\begin{tikzpicture}[x=1cm,y=1cm]
  \path[use as bounding box] (0,-0.88) rectangle (12.75,4.12);

  \foreach \x/\name/\lab in {1.15/Cone/C_1,3.55/Ctwo/C_2,5.95/Cthree/C_3,8.65/Cj/C_j,11.35/Clast/C_{k^2}}{
    \node[draw=ThLine,rounded corners=1pt,minimum width=1.28cm,
      minimum height=2.08cm,fill=ThMist] (\name) at (\x,0.42) {};
    \node[below=1.2mm of \name,font=\fontsize{9pt}{11pt}\selectfont\bfseries,
      text=ThInkSoft] {$\mathbf{\lab}$};
  }
  \node[font=\fontsize{13pt}{15pt}\selectfont,text=ThMuted] at (7.30,0.42) {$\cdots$};
  \node[font=\fontsize{13pt}{15pt}\selectfont,text=ThMuted] at (10.00,0.42) {$\cdots$};

  \node[thStoch,minimum size=0.36cm] at ($(Cone.center)+(0,0.00)$) {};
  \node[thStoch,minimum size=0.36cm] at ($(Cone.center)+(0,-0.60)$) {};
  \node[thStoch,minimum size=0.36cm] at (Ctwo.center) {};
  \foreach \dy in {0.60,0.00,-0.60}{
    \node[thStoch,minimum size=0.36cm] at ($(Cthree.center)+(0,\dy)$) {};
  }
  \node[thStoch,minimum size=0.36cm] at ($(Cj.center)+(0,0.60)$) {};
  \node[thStoch,minimum size=0.36cm] at ($(Cj.center)+(0,-0.60)$) {};
  \node[thStoch,minimum size=0.36cm] at ($(Clast.center)+(0,0.60)$) {};
  \node[thStoch,minimum size=0.36cm] at (Clast.center) {};

  \foreach \name in {Cone,Ctwo,Cj,Clast}{
    \node[draw=ThBlue,line width=1.15pt,rounded corners=2pt,
      inner sep=1.8mm,fit=(\name)] {};
  }

  \node[thPacket,minimum size=0.50cm] (pOne) at ($(Cone.center)+(0,0.60)$) {$z_1$};
  \node[thPacket,minimum size=0.50cm] (pTwo) at ($(Ctwo.center)+(0,0.60)$) {$z_2$};
  \node[thPacket,minimum size=0.50cm] (pThree) at ($(Ctwo.center)+(0,-0.60)$) {$z_3$};
  \node[thPacket,minimum size=0.50cm] (pKmOne) at (Cj.center) {};
  \node[font=\fontsize{5.6pt}{6.4pt}\selectfont,text=ThWhite,inner sep=0pt]
    at (pKmOne.center) {$z_{\!k-1}$};
  \node[thPacket,minimum size=0.50cm] (pK) at ($(Clast.center)+(0,-0.60)$) {$z_k$};

  \node[thPacket,minimum size=0.58cm] (zOne) at (1.15,3.12) {$z_1$};
  \node[thPacket,minimum size=0.58cm] (zTwo) at (3.05,3.12) {$z_2$};
  \node[thPacket,minimum size=0.58cm] (zThree) at (4.55,3.12) {$z_3$};
  \node[font=\fontsize{12pt}{14pt}\selectfont,text=ThGold] at (6.35,3.12) {$\cdots$};
  \node[thPacket,minimum size=0.58cm] (zKmOne) at (8.25,3.12) {};
  \node[font=\fontsize{6.6pt}{7.4pt}\selectfont,text=ThWhite,inner sep=0pt]
    at (zKmOne.center) {$z_{\!k-1}$};
  \node[thPacket,minimum size=0.58cm] (zK) at (10.85,3.12) {$z_k$};

  \draw[thArrowThin,draw=ThGold,line width=0.90pt] (zOne)--(zTwo);
  \draw[thArrowThin,draw=ThGold,line width=0.90pt] (zTwo)--(zThree);
  \draw[draw=ThGold,line width=0.90pt] (zThree.east)--(5.62,3.12);
  \draw[thArrowThin,draw=ThGold,line width=0.90pt] (7.08,3.12)--(zKmOne.west);
  \draw[thArrowThin,draw=ThGold,line width=0.90pt] (zKmOne)--(zK);

  \node[draw=ThGold,line width=0.9pt,rounded corners=2pt,
    inner xsep=3.6mm,inner ysep=2.6mm,fit=(zOne)(zK)] (packetBox) {};
  \node[above=1.2mm of packetBox,font=\fontsize{10pt}{12pt}\selectfont\bfseries,
    text=ThGold] {one packet \(\mathbf{P}\)};
\end{tikzpicture}}

\newcommand{\ThManyRepresentedCase}{%
\begin{tikzpicture}[x=1cm,y=1cm]
  \path[use as bounding box] (0,-0.48) rectangle (5.95,1.25);
  \node[thPacket,minimum size=0.58cm] (ya) at (0.62,0.30) {$y_a$};
  \node[thPacket,minimum size=0.43cm] (ybOne) at (2.72,0.88) {};
  \node[thPacket,minimum size=0.43cm] (ybTwo) at (3.72,0.53) {};
  \node[thPacket,minimum size=0.43cm] (ybLast) at (5.18,-0.02) {};
  \foreach \target in {ybOne,ybTwo,ybLast}{
    \draw[thArrowThin,draw=ThBlue] (ya)--(\target);
  }
  \node[anchor=west,font=\fontsize{8.3pt}{9.8pt}\selectfont,text=ThInkSoft]
    at ([xshift=1.2mm]ybOne.east) {$y_{a+1}$};
  \node[anchor=west,font=\fontsize{8.3pt}{9.8pt}\selectfont,text=ThInkSoft]
    at ([xshift=1.2mm]ybTwo.east) {$y_{a+2}$};
  \node[font=\fontsize{10pt}{12pt}\selectfont,text=ThMuted] at (4.63,0.18) {$\vdots$};
  \node[anchor=west,font=\fontsize{8.3pt}{9.8pt}\selectfont,text=ThInkSoft]
    at ([xshift=1.2mm]ybLast.east) {$y_s$};
\end{tikzpicture}}

\newcommand{\ThRepeatedColumnsCase}{%
\begin{tikzpicture}[x=1cm,y=1cm]
  \path[use as bounding box] (0,-0.48) rectangle (5.95,1.25);
  \foreach \x/\name in {0.72/repOne,2.18/repTwo,3.64/repThree,5.10/repFour}{
    \node[draw=ThBlue,rounded corners=1pt,minimum width=0.80cm,
      minimum height=1.20cm,fill=ThMist] (\name) at (\x,0.30) {};
  }
  \node[thPacket,minimum size=0.29cm] (repOneTop) at ($(repOne.center)+(0,0.37)$) {};
  \node[thPacket,minimum size=0.29cm] at (repOne.center) {};
  \node[thPacket,minimum size=0.29cm] (repOneBottom) at ($(repOne.center)+(0,-0.37)$) {};
  \node[thPacket,minimum size=0.31cm] at ($(repTwo.center)+(0,0.27)$) {};
  \node[thPacket,minimum size=0.31cm] at ($(repTwo.center)+(0,-0.27)$) {};
  \node[thPacket,minimum size=0.31cm] at (repThree.center) {};
  \node[thPacket,minimum size=0.31cm] at ($(repFour.center)+(0,0.27)$) {};
  \node[thPacket,minimum size=0.31cm] at ($(repFour.center)+(0,-0.27)$) {};
  \draw[-{Stealth[length=1.7mm]},draw=ThCoral,line width=0.82pt]
    ([xshift=0.7mm]repOneTop.east) to[bend left=48]
    node[midway,right=0.5mm,font=\fontsize{8.3pt}{9.8pt}\selectfont,
      text=ThCoral] {$k^2$}
    ([xshift=0.7mm]repOneBottom.east);
\end{tikzpicture}}

\newcommand{\ThPacketShortcutEvent}{%
\begin{tikzpicture}[x=1cm,y=1cm]
  \path[use as bounding box] (0,-0.72) rectangle (9.55,1.65);
  \node[thPacketActive,minimum size=0.58cm]
    (za) at (1.05,0.22) {$z_a$};
  \node[thInactive,draw=ThGold!78!ThInk,line width=0.85pt,minimum size=0.48cm]
    (midOne) at (3.00,0.22) {$\times$};
  \node[thInactive,draw=ThGold!78!ThInk,line width=0.85pt,minimum size=0.48cm]
    (midTwo) at (4.15,0.22) {$\times$};
  \node[font=\fontsize{11pt}{13pt}\selectfont,text=ThMuted] at (5.15,0.22) {$\cdots$};
  \node[thInactive,draw=ThGold!78!ThInk,line width=0.85pt,minimum size=0.48cm]
    (midLast) at (6.15,0.22) {$\times$};
  \node[thPacketActive,minimum size=0.58cm]
    (zb) at (8.35,0.22) {$z_b$};

  \draw[ThLine,line width=0.72pt] (za)--(midOne)--(midTwo);
  \draw[ThLine,line width=0.72pt] (midTwo.east)--(4.63,0.22);
  \draw[ThLine,line width=0.72pt] (5.67,0.22)--(midLast.west);
  \draw[-{Stealth[length=1.7mm]},draw=ThLine,line width=0.72pt] (midLast)--(zb);
  \draw[-{Stealth[length=2.0mm]},draw=ThBlue,line width=1.05pt]
    (za.north) to[bend left=24]
    node[midway,above=1.0mm,font=\fontsize{9pt}{10.8pt}\selectfont\bfseries,
      text=ThBlue] {realized shortcut}
    (zb.north);

  \node[below=1.2mm of za,font=\fontsize{8.2pt}{9.7pt}\selectfont,
    text=ThTeal] {active};
  \node[below=1.2mm of zb,font=\fontsize{8.2pt}{9.7pt}\selectfont,
    text=ThTeal] {active};
  \draw[decorate,decoration={brace,mirror,amplitude=3.0pt},draw=ThMuted,
    line width=0.65pt]
    ([yshift=-2.4mm]midOne.south west)--
    node[midway,below=1.3mm,font=\fontsize{8.2pt}{9.7pt}\selectfont,
      text=ThMuted] {all inactive}
    ([yshift=-2.4mm]midLast.south east);
\end{tikzpicture}}

% --------------------------------------------------------------------------
% Layered-column stage.
% --------------------------------------------------------------------------
\newcommand{\ThLayeredStage}{%
\begin{tikzpicture}[>=Stealth,x=1cm,y=1cm]
  \node[thDet] (L1) at (0,1.05) {$X_{i-1}^1$};
  \node[thDet] (L2) at (0,-1.05) {$X_{i-1}^2$};
  \node[thStoch,minimum size=0.63cm] (A) at (2.45,1.8) {$1,*$};
  \node[thStoch,minimum size=0.63cm] (C) at (2.45,0.6) {$*,1$};
  \node[thStoch,minimum size=0.63cm] (D) at (2.45,-0.6) {$*,2$};
  \node[thStoch,minimum size=0.63cm] (B) at (2.45,-1.8) {$2,*$};
  \node[thDet] (R1) at (4.90,1.05) {$X_i^1$};
  \node[thDet] (R2) at (4.90,-1.05) {$X_i^2$};
  \foreach \u/\v in {L1/A,L1/C,L1/D,L2/B,L2/C,L2/D,A/R1,A/R2,B/R1,B/R2,C/R1,D/R2}{
    \draw[-{Stealth[length=1.5mm]},ThMuted,line width=0.65pt] (\u)--(\v);
  }
  \node[thSmallLabel] at (0,-2.35) {gate column};
  \node[thSmallLabel] at (2.45,-2.35) {routing column};
  \node[thSmallLabel] at (4.90,-2.35) {gate column};
\end{tikzpicture}}

% --------------------------------------------------------------------------
% Asymmetric detours and circulation pipeline.
% --------------------------------------------------------------------------
\newcommand{\ThSameAnchorAttachment}{%
\begin{tikzpicture}[x=1cm,y=1cm]
  \path[use as bounding box] (0,-0.35) rectangle (5,1.55);
  \node[thDet] (sa) at (1.1,0.65) {$a$};
  \node[thStoch] (si) at (3.2,0.65) {$i$};
  \draw[thArrow,draw=ThCoral] ([yshift=0.12cm]sa.east)
    to[bend left=22] node[above,thSmallLabel] {$d(a,i)$} ([yshift=0.12cm]si.west);
  \draw[thArrow,draw=ThCoral] ([yshift=-0.12cm]si.west)
    to[bend left=22] node[below,thSmallLabel] {$d(i,a)$} ([yshift=-0.12cm]sa.east);
\end{tikzpicture}}

\newcommand{\ThTwoEndpointDetour}{%
\begin{tikzpicture}[x=1cm,y=1cm]
  \path[use as bounding box] (0,-0.35) rectangle (5,1.55);
  \node[thDet] (a) at (0.50,0.65) {$a$};
  \node[thStoch] (i) at (2.50,0.65) {$i$};
  \node[thDet] (b) at (4.50,0.65) {$b$};
  \draw[thArrow,draw=ThTeal] (a)--(i);
  \draw[thArrow,draw=ThTeal] (i)--(b);
  \draw[thArrowThin,draw=ThBlue,dashed] ([yshift=-0.18cm]a.south)
    to[bend right=20]
    ([yshift=-0.28cm]b.south);
\end{tikzpicture}}

\newcommand{\ThCirculationPipeline}{%
\begin{tikzpicture}[x=1cm,y=1cm]
  \tikzset{pipe/.style={rounded corners=1pt,draw=ThLine,fill=ThWhite,
    minimum width=1.72cm,minimum height=0.90cm,align=center,
    font=\fontsize{9.4pt}{11pt}\selectfont\bfseries}}
  \node[pipe,draw=ThTeal!55,fill=ThTeal!7] (sample) at (0.90,0.65)
    {sample a\\backbone};
  \node[pipe,draw=ThBlue!55,fill=ThBlue!7] (flow) at (3.20,0.65)
    {fractional\\circulation};
  \node[pipe,draw=ThBlue!55,fill=ThBlue!7] (couple) at (5.50,0.65)
    {assign detour\\endpoints};
  \node[pipe,draw=ThGold!60,fill=ThGold!9] (round) at (7.80,0.65)
    {joint\\rounding};
  \node[pipe,draw=ThCoral!55,fill=ThCoral!7] (probe) at (10.10,0.65)
    {adaptive\\probing};
  \foreach \u/\v in {sample/flow,flow/couple,couple/round,round/probe}{
    \draw[thArrowThin] (\u)--(\v);
  }
  \node[thSmallLabel,text=ThTeal] at (0.90,-0.32) {backbone first};
  \node[thSmallLabel,text=ThBlue] at (5.50,-0.32) {two-endpoint detours};
  \node[thSmallLabel,text=ThGold!78!ThInk] at (7.80,-0.32)
    {interval control};
  \node[thSmallLabel,text=ThCoral] at (10.10,-0.32) {activations last};
\end{tikzpicture}}

% --------------------------------------------------------------------------
% Asymmetric clairvoyance upper bound: circulation and rounding details.
% --------------------------------------------------------------------------
\newcommand{\ThSampledBackboneCycle}{%
\begin{tikzpicture}[x=1cm,y=1cm]
  \path[use as bounding box] (-2.45,-2.30) rectangle (2.45,2.30);

  \coordinate (r)  at (180:1.86);
  \coordinate (s1) at (135:1.86);
  \coordinate (s2) at ( 90:1.86);
  \coordinate (s3) at ( 45:1.86);
  \coordinate (s4) at (  0:1.86);
  \coordinate (s5) at (-45:1.86);
  \coordinate (s6) at (-90:1.86);
  \coordinate (s7) at (-135:1.86);

  \begin{scope}[on background layer]
    \draw[ThBlue,line width=1.22pt] (0,0) circle[radius=1.86cm];
    \draw[thArrowThin,draw=ThBlue,line width=1.22pt]
      (116:1.86) arc[start angle=116,end angle=104,radius=1.86cm];
    \draw[thArrowThin,draw=ThBlue,line width=1.22pt]
      (-64:1.86) arc[start angle=-64,end angle=-76,radius=1.86cm];
  \end{scope}

  \node[thDepot,minimum size=0.50cm] at (r) {$r$};
  \node[thDet] at (s1) {$s_1$};
  \node[thDet] at (s2) {$s_2$};
  \node[thDet] at (s3) {$s_3$};
  \node[thDet] at (s4) {$s_4$};
  \node[thDet] at (s5) {$s_5$};
  \node[thDet] at (s6) {$s_6$};
  \node[thDet] at (s7) {$s_7$};
  \node[font=\fontsize{12pt}{14pt}\selectfont\bfseries,text=ThBlue]
    at (0,0) {optimal tour \(\mathbf{T_R}\)};
\end{tikzpicture}}

\newcommand{\ThCirculationRelaxation}{%
\begin{tikzpicture}[x=1cm,y=1cm]
  \path[use as bounding box] (-2.45,-2.30) rectangle (2.45,2.30);

  \coordinate (a) at (135:1.86);
  \coordinate (e) at ( 90:1.86);
  \coordinate (b) at ( 45:1.86);
  \coordinate (f) at (  0:1.86);
  \coordinate (c) at (-45:1.86);
  \coordinate (d) at (-90:1.86);
  \coordinate (r) at (-135:1.86);
  \coordinate (g) at (180:1.86);
  \coordinate (i) at (0,0);

  \begin{scope}[on background layer]
    \draw[ThLine,line width=0.76pt] (0,0) circle[radius=1.86cm];
    \draw[thArrowThin,draw=ThLine,line width=0.76pt]
      (116:1.86) arc[start angle=116,end angle=104,radius=1.86cm];
  \end{scope}

  \node[thDet] (na) at (a) {$a$};
  \node[thDet] at (e) {$e$};
  \node[thDet] (nb) at (b) {$b$};
  \node[thDet] at (f) {$f$};
  \node[thDet] (nc) at (c) {$c$};
  \node[thDet] at (d) {$d$};
  \node[thDepot] (nr) at (r) {$r$};
  \node[thDet] at (g) {$g$};
  \node[thStoch,minimum size=0.60cm] (ni) at (i) {$i$};

  \draw[thArrow,draw=ThBlue,line width=1.05pt]
    (na.south east) to[bend right=7]
    node[pos=0.42,above=1pt,thSmallLabel,text=ThBlue] {$z_{ai}$}
    (ni.north west);
  \draw[thArrow,draw=ThBlue,line width=1.05pt]
    (nc.north west) to[bend right=7]
    node[pos=0.42,below=1pt,thSmallLabel,text=ThBlue] {$z_{ci}$}
    (ni.south east);
  \draw[thArrow,draw=ThTeal,line width=1.05pt]
    (ni.north east) to[bend right=7]
    node[pos=0.58,above=1pt,thSmallLabel,text=ThTeal] {$z_{ib}$}
    (nb.south west);
  \draw[thArrow,draw=ThTeal,line width=1.05pt]
    (ni.south west) to[bend right=7]
    node[pos=0.58,below=1pt,thSmallLabel,text=ThTeal] {$z_{ir}$}
    (nr.north east);

  \node[thSmallLabel,text=ThBlue,anchor=east] at (-0.38,0.10)
    {incoming};
  \node[thSmallLabel,text=ThTeal,anchor=west] at (0.38,0.10)
    {outgoing};
\end{tikzpicture}}

\newcommand{\ThLongestRunWitness}{%
\begin{tikzpicture}[x=1cm,y=1cm]
  \path[use as bounding box] (0,-1.42) rectangle (5.15,1.14);
  \coordinate (s0) at (0.38,0.52);
  \coordinate (u1) at (1.12,0.52);
  \coordinate (u2) at (1.84,0.52);
  \coordinate (u3) at (2.56,0.52);
  \coordinate (s1) at (3.30,0.52);
  \coordinate (v1) at (4.02,0.52);
  \coordinate (s2) at (4.76,0.52);
  \begin{scope}[on background layer]
    \draw[thArrowThin,draw=ThBlue,line width=0.90pt]
      (s0)--(u1)--(u2)--(u3)--(s1)--(v1)--(s2);
    \draw[thArrowThin,draw=ThGold!80!ThInk,densely dashed,line width=0.78pt]
      (s0) to[bend left=20] (s1);
    \draw[thArrowThin,draw=ThGold!80!ThInk,densely dashed,line width=0.78pt]
      (s1) to[bend left=20] (s2);
  \end{scope}
  \node[thDet,minimum width=0.48cm,minimum height=0.44cm] at (s0) {$s_0$};
  \node[thStoch,minimum size=0.42cm] at (u1) {$u_1$};
  \node[thStoch,minimum size=0.42cm] at (u2) {$u_2$};
  \node[thStoch,minimum size=0.42cm] at (u3) {$u_3$};
  \node[thDet,minimum width=0.48cm,minimum height=0.44cm] at (s1) {$s_1$};
  \node[thStoch,minimum size=0.42cm] at (v1) {$v_1$};
  \node[thDet,minimum width=0.48cm,minimum height=0.44cm] at (s2) {$s_2$};
  \draw[decorate,decoration={brace,amplitude=4pt,mirror},ThTeal,line width=0.85pt]
    (0.72,0.08)--(2.96,0.08)
    node[midway,below=0.18cm,thSmallLabel,text=ThTeal] {run of length 3};
  \draw[decorate,decoration={brace,amplitude=4pt,mirror},ThMuted,line width=0.75pt]
    (3.62,0.08)--(4.42,0.08)
    node[midway,below=0.18cm,thSmallLabel] {run of length 1};
  \node[thSmallLabel,text=ThGold!80!ThInk] at (2.55,-1.12)
    {dashed: backbone shortcuts};
\end{tikzpicture}}

\newcommand{\ThTransformArrow}{%
\begin{tikzpicture}[x=1cm,y=1cm]
  \path[use as bounding box] (0,-0.20) rectangle (1.05,0.20);
  \draw[thArrow,draw=ThInkSoft,line width=0.95pt] (0.06,0)--(0.99,0);
\end{tikzpicture}}

\newcommand{\ThParallelDetours}{%
\begin{tikzpicture}[x=1cm,y=1cm]
  \path[use as bounding box] (0,-1.10) rectangle (3.70,1.10);

  \node[thDet,minimum width=0.48cm,minimum height=0.44cm] (ls) at (0.32,0)
    {$s_0$};
  \node[thStoch,minimum size=0.42cm] (u1) at (1.85,0.82) {$u_1$};
  \node[thStoch,minimum size=0.42cm] (u2) at (1.85,0) {$u_2$};
  \node[thStoch,minimum size=0.42cm] (u3) at (1.85,-0.82) {$u_3$};
  \node[thDet,minimum width=0.48cm,minimum height=0.44cm] (rs) at (3.38,0)
    {$s_1$};

  \foreach \u in {u1,u2,u3}{
    \draw[thArrowThin,draw=ThBlue,line width=0.88pt]
      (ls.east)--(\u.west);
    \draw[thArrowThin,draw=ThTeal,line width=0.88pt]
      (\u.east)--(rs.west);
  }
\end{tikzpicture}}

\newcommand{\ThSegmentPadding}{%
\begin{tikzpicture}[x=1cm,y=1cm]
  \path[use as bounding box] (0,-1.65) rectangle (5.25,1.28);
  \node[thDet] (ls) at (0.38,0.18) {$s_j$};
  \node[thStoch] (lu) at (1.86,0.92) {$u_1$};
  \node[thStoch] (lv) at (2.98,-0.62) {$u_2$};
  \node[thDet] (rs) at (4.82,0.18) {$s_{j+1}$};
  \begin{scope}[on background layer]
    \draw[thArrowThin,draw=ThTeal,line width=1.00pt] (ls)--(lu)--(rs);
    \draw[thArrowThin,draw=ThBlue,line width=1.00pt] (ls)--(lv)--(rs);
    \draw[thArrowThin,draw=ThGold!80!ThInk,dashed,line width=0.90pt]
      ([yshift=-0.14cm]ls.south)
      .. controls (1.30,-1.55) and (3.90,-1.55)
      .. ([yshift=-0.14cm]rs.south)
      node[midway,below=0.08cm,thSmallLabel,text=ThGold!80!ThInk]
      {\(M-k_j\) direct flow};
  \end{scope}
  \node[thLabel,anchor=west] at ([xshift=1.8mm]lu.north east)
    {\(k_j=2,\ M=3\)};
\end{tikzpicture}}

% Unit-flow detours between consecutive vertices of R in T_C order.
% The segment counts are (2,3,1); distinct direct copies pad every total to 3.
\newcommand{\ThPaddedCirculationCycle}{%
\begin{tikzpicture}[x=1cm,y=1cm,
  thCycleOutside/.style={thStoch,minimum size=0.28cm,inner sep=0pt}]
  \path[use as bounding box] (-3.35,-1.80) rectangle (3.35,2.30);
  \node[thDepot,minimum width=0.88cm] (s0) at (-2.65,1.15) {$s_0=r$};
  \node[thDet] (s1) at (2.65,1.15) {$s_1$};
  \node[thDet] (s2) at (0,-1.40) {$s_2$};

  \node[thCycleOutside] (u01) at (0,1.95) {};
  \node[thCycleOutside] (u02) at (0,0.58) {};
  \node[thCycleOutside] (u11) at (0.90,0) {};
  \node[thCycleOutside] (u12) at (1.50,-0.45) {};
  \node[thCycleOutside] (u13) at (2.12,-0.90) {};
  \node[thCycleOutside] (u21) at (-1.30,-0.16) {};

  \begin{scope}[on background layer]
    \foreach \u/\shade in {u01/ThTeal,u02/ThBlue}{
      \draw[thArrowThin,draw=\shade,line width=0.90pt] (s0)--(\u);
      \draw[thArrowThin,draw=\shade,line width=0.90pt] (\u)--(s1);
    }
    \foreach \u/\shade in {u11/ThTeal,u12/ThBlue,u13/ThTeal}{
      \draw[thArrowThin,draw=\shade,line width=0.90pt] (s1)--(\u);
      \draw[thArrowThin,draw=\shade,line width=0.90pt] (\u)--(s2);
    }
    \draw[thArrowThin,draw=ThBlue,line width=0.90pt] (s2)--(u21);
    \draw[thArrowThin,draw=ThBlue,line width=0.90pt] (u21)--(s0);

    \draw[thArrowThin,draw=ThGold!80!ThInk,dashed,line width=0.90pt]
      (s0.east)--(s1.west);
    \draw[thArrowThin,draw=ThGold!80!ThInk,dashed,line width=0.90pt]
      ([yshift=0.08cm]s2.west) to[bend left=16]
      ([xshift=0.11cm]s0.south);
    \draw[thArrowThin,draw=ThGold!80!ThInk,dashed,line width=0.90pt]
      ([yshift=-0.08cm]s2.west) to[bend left=33]
      ([xshift=-0.13cm]s0.south);
  \end{scope}
\end{tikzpicture}}

\newcommand{\ThJointRounding}{%
\begin{tikzpicture}[x=1cm,y=1cm]
  \path[use as bounding box] (0,-0.96) rectangle (11,1.48);

  \node[thDet] (la) at (0.72,0.64) {$a$};
  \node[thDet] (lb) at (2.35,1.02) {$b$};
  \node[thDet] (lc) at (3.98,0.64) {$c$};
  \node[thStoch] (li) at (2.35,-0.64) {$i$};
  \begin{scope}[on background layer]
    \draw[thArrowThin,draw=ThBlue!62,line width=0.86pt]
      (la) to[bend right=8] (li);
    \draw[thArrowThin,draw=ThBlue!62,line width=0.86pt]
      (li) to[bend right=8] (lb);
    \draw[thArrowThin,draw=ThTeal!68,line width=0.86pt]
      (lc) to[bend left=8] (li);
    \draw[thArrowThin,draw=ThTeal!68,line width=0.86pt]
      (li) to[bend left=20] (la);
    \draw[thArrowThin,draw=ThGold!70!ThInk,dashed,line width=0.76pt]
      (lc) to[bend right=25]
      node[midway,above=2mm,thSmallLabel,text=ThGold!75!ThInk] {$y_{ca}$} (la);
  \end{scope}

  \draw[thArrow,draw=ThInkSoft,line width=1.0pt] (4.68,0.16)--(6.18,0.16);

  \node[thDet] (ra) at (7.02,0.64) {$a_i$};
  \node[thDet] (rb) at (8.65,1.02) {$b_i$};
  \node[thDet] (rc) at (10.28,0.64) {$c$};
  \node[thStoch] (ri) at (8.65,-0.64) {$i$};
  \begin{scope}[on background layer]
    \draw[thArrow,draw=ThTeal,line width=1.15pt] (ra)--(ri)--(rb);
    \draw[thArrowThin,draw=ThBlue,dashed,line width=0.82pt]
      (ra) to[bend left=18]
      node[midway,above,thSmallLabel,text=ThBlue] {collapsed chord} (rb);
    \draw[thArrowThin,draw=ThGold!80!ThInk,line width=0.90pt]
      (rc.south) to[bend left=24]
      node[pos=0.18,below=0.04cm,thSmallLabel,text=ThGold!80!ThInk]
      {\(N_{ca}\)} (ra.south);
  \end{scope}
\end{tikzpicture}}

\newcommand{\ThDyadicRounding}{%
\begin{tikzpicture}[x=1cm,y=1cm]
  \path[use as bounding box] (0,-1.72) rectangle (11,1.55);
  \foreach \j in {0,...,7}{
    \coordinate (s\j) at ({1.18+1.23*\j},0.28);
  }
  \draw[ThLine,line width=0.72pt] (s0)--(s7);
  \foreach \j in {0,...,7}{
    \node[thDet,minimum width=0.43cm,minimum height=0.40cm,
      font=\fontsize{7.6pt}{8.5pt}\selectfont] at (s\j) {$s_{\j}$};
  }

  \draw[decorate,decoration={brace,amplitude=5pt},ThGold!82!ThInk,line width=0.95pt]
    ([yshift=0.28cm]s1.north)--([yshift=0.28cm]s6.north)
    node[midway,above=0.18cm,thLabel,text=ThGold!82!ThInk] {one cyclic interval \(I\)};

  \foreach \a/\b in {0/1,2/3,4/5,6/7}{
    \draw[decorate,decoration={brace,amplitude=3.5pt,mirror},ThTeal,line width=0.68pt]
      ([yshift=-0.22cm]s\a.south)--([yshift=-0.22cm]s\b.south);
  }
  \foreach \a/\b in {0/3,4/7}{
    \draw[decorate,decoration={brace,amplitude=4pt,mirror},ThBlue,line width=0.75pt]
      ([yshift=-0.70cm]s\a.south)--([yshift=-0.70cm]s\b.south);
  }
  \draw[decorate,decoration={brace,amplitude=4.5pt,mirror},ThMuted,line width=0.78pt]
    ([yshift=-1.18cm]s0.south)--([yshift=-1.18cm]s7.south)
    node[midway,below=0.14cm,thSmallLabel] {dyadic interval family};
\end{tikzpicture}}

\newcommand{\ThIntervalImbalance}{%
\begin{tikzpicture}[x=1cm,y=1cm]
  \path[use as bounding box] (0,-0.72) rectangle (11,1.58);
  \foreach \j in {0,...,6}{
    \coordinate (s\j) at ({0.82+1.56*\j},0.05);
  }
  \begin{scope}[on background layer]
    \draw[thArrowThin,draw=ThLine,line width=0.72pt]
      (s0)--(s1)--(s2)--(s3)--(s4)--(s5)--(s6);
    \draw[thArrow,draw=ThTeal,line width=1.05pt]
      (s0) to[bend left=25]
      node[midway,above,thSmallLabel,text=ThTeal]
      {entering: \(\chi_I=-1\)} (s3);
    \draw[thArrow,draw=ThCoral,line width=1.05pt]
      (s3) to[bend left=25]
      node[midway,above,thSmallLabel,text=ThCoral]
      {leaving: \(\chi_I=+1\)} (s6);
  \end{scope}
  \foreach \j in {0,1,5,6}{
    \node[thDet,minimum width=0.43cm,minimum height=0.40cm] at (s\j) {$s_{\j}$};
  }
  \foreach \j in {2,3,4}{
    \node[thDet,fill=ThTeal!12,minimum width=0.43cm,minimum height=0.40cm]
      at (s\j) {$s_{\j}$};
  }
  \draw[decorate,decoration={brace,amplitude=5pt},ThTeal,line width=0.88pt]
    ([yshift=0.24cm]s2.north)--([yshift=0.24cm]s4.north)
    node[midway,above=0.17cm,thLabel,text=ThTeal] {cyclic interval \(I\)};
\end{tikzpicture}}

\newcommand{\ThAdaptiveChordWalk}{%
\begin{tikzpicture}[x=1cm,y=1cm]
  \path[use as bounding box] (0,-0.30) rectangle (11,1.52);
  \draw[ThLine,line width=0.65pt] (5.50,-0.24)--(5.50,1.46);

  \node[thDet] (la) at (0.85,0.10) {$a$};
  \node[thActive] (li) at (2.63,0.96) {$i$};
  \node[thDet] (lb) at (4.40,0.10) {$b_i$};
  \begin{scope}[on background layer]
    \draw[thArrow,draw=ThTeal,line width=1.38pt] (la)--(li)--(lb);
  \end{scope}
  \fill[ThGold!78!ThInk]
    (4.40,0.63)--(4.23,0.94)--(4.57,0.94)--cycle;
  \node[thSmallLabel,text=ThGold!75!ThInk] at (4.40,1.18) {current position};

  \node[thDet] (ra) at (6.60,0.10) {$a$};
  \node[thInactive] (ri) at (8.37,0.96) {$i$};
  \node[thDet] (rb) at (10.15,0.10) {$b_i$};
  \fill[ThGold!78!ThInk]
    (6.60,0.63)--(6.43,0.94)--(6.77,0.94)--cycle;
  \node[thSmallLabel,text=ThGold!75!ThInk] at (6.60,1.18) {current position};
\end{tikzpicture}}

% Initial state on overlay 2, then a complete virtual walk on overlays 3--17.
\newcommand{\ThAuxiliaryPointerSteps}{%
\begin{tikzpicture}[x=1cm,y=1cm]
  \path[use as bounding box] (0,-2.20) rectangle (11,3.03);

  \node[thDepot,minimum width=0.85cm,minimum height=0.52cm] (pa) at (1,0) {$a=r$};
  \node[thDet,minimum width=0.60cm,minimum height=0.52cm] (pnext) at (5.5,0) {$a^+$};
  \node[thDet,minimum width=0.60cm,minimum height=0.52cm] (pb) at (10,0) {$b$};
  \node[thDet,minimum width=0.60cm,minimum height=0.52cm] (purone) at (8.5,0.633) {$b^+$};
  \node[thDet,minimum width=0.60cm,minimum height=0.52cm] (pultwo) at (2.5,0.633) {$a^-$};
  \node[thDet,minimum width=0.70cm,minimum height=0.44cm] (pbrone) at (7,0) {$a^{++}$};
  \node[thDet,minimum width=0.65cm,minimum height=0.44cm] (pbrtwo) at (8.5,0) {$b^-$};

  % Keep the same backbone vertices and positions on every overlay.
  \foreach \name/\x/\y in {
    purtwo/7/1.267,pr/5.5/1.90,pulone/4/1.267}{
    \node[thDet,minimum width=0.40cm,minimum height=0.40cm,inner sep=0pt]
      (\name) at (\x,\y) {};
  }

  % Outside vertices i and j are active; k is inactive. Fading means probed.
  \only<2>{
    \node[thStoch,minimum size=0.40cm,inner sep=0pt] (pi) at (3,-0.52) {$i$};
  }
  \only<3>{
    \node[thActive,minimum size=0.40cm,inner sep=0pt] (pi) at (3,-0.52) {$i$};
  }
  \only<4-17>{
    \node[thStoch,draw=ThMuted!60,fill=ThMuted!12,text=ThMuted,
      minimum size=0.40cm,inner sep=0pt] (pi) at (3,-0.52) {$i$};
  }
  \only<2-8>{
    \node[thStoch,minimum size=0.40cm,inner sep=0pt] (pj) at (5.5,2.68) {$j$};
  }
  \only<9>{
    \node[thActive,minimum size=0.40cm,inner sep=0pt] (pj) at (5.5,2.68) {$j$};
  }
  \only<10-17>{
    \node[thStoch,draw=ThMuted!60,fill=ThMuted!12,text=ThMuted,
      minimum size=0.40cm,inner sep=0pt] (pj) at (5.5,2.68) {$j$};
  }
  \only<2-4>{
    \node[thStoch,minimum size=0.40cm,inner sep=0pt] (pk) at (7.75,-0.65) {};
    \node[thSmallLabel,anchor=west,text=ThInk] at (8.02,-0.65) {$k$};
  }
  \only<5>{
    \node[thInactive,draw=ThCoral,fill=ThCoral!10,text=ThCoral,
      minimum size=0.40cm,inner sep=0pt] (pk) at (7.75,-0.65) {$\times$};
    \node[thSmallLabel,anchor=west,text=ThCoral] at (8.02,-0.65) {$k$};
  }
  \only<6-17>{
    \node[thInactive,draw=ThMuted!60,fill=ThMuted!12,text=ThMuted,
      minimum size=0.40cm,inner sep=0pt] (pk) at (7.75,-0.65) {$\times$};
    \node[thSmallLabel,anchor=west,text=ThMuted] at (8.02,-0.65) {$k$};
  }

  \begin{scope}[on background layer]
    \foreach \u/\v in {pa/pnext,pnext/pbrone,pbrone/pbrtwo,pbrtwo/pb,
      pb/purone,purone/purtwo,purtwo/pr,pr/pulone,pulone/pultwo,pultwo/pa}{
      \draw[thArrowThin,draw=ThBlue!50!ThMuted,line width=0.85pt] (\u)--(\v);
    }
    \only<2>{
      \foreach \u/\v in {pa/pi,pi/pnext}{
        \draw[thArrowThin,draw=ThTeal!75,line width=0.75pt] (\u)--(\v);
      }
    }
    \only<3>{
      \foreach \u/\v in {pa/pi,pi/pnext}{
        \draw[thArrow,draw=ThTeal,line width=2.0pt] (\u)--(\v);
      }
    }
    \only<4-17>{
      \foreach \u/\v in {pa/pi,pi/pnext}{
        \draw[thArrowThin,draw=ThMuted!45,densely dashed,line width=0.70pt]
          (\u)--(\v);
      }
    }
    \only<2-8>{
      \foreach \u/\v in {purone/pj,pj/pultwo}{
        \draw[thArrowThin,draw=ThTeal!75,line width=0.75pt] (\u)--(\v);
      }
    }
    \only<9>{
      \foreach \u/\v in {purone/pj,pj/pultwo}{
        \draw[thArrow,draw=ThTeal,line width=2.0pt] (\u)--(\v);
      }
    }
    \only<10-17>{
      \foreach \u/\v in {purone/pj,pj/pultwo}{
        \draw[thArrowThin,draw=ThMuted!45,densely dashed,line width=0.70pt]
          (\u)--(\v);
      }
    }
    \only<2-4>{
      \foreach \u/\v in {pbrone/pk,pk/pbrtwo}{
        \draw[thArrowThin,draw=ThTeal!75,line width=0.75pt] (\u)--(\v);
      }
    }
    \only<5>{
      \foreach \u/\v in {pbrone/pk,pk/pbrtwo}{
        \draw[thArrowThin,draw=ThCoral!70,densely dashed,line width=0.70pt]
          (\u)--(\v);
      }
    }
    \only<6-17>{
      \foreach \u/\v in {pbrone/pk,pk/pbrtwo}{
        \draw[thArrowThin,draw=ThMuted!45,densely dashed,line width=0.70pt]
          (\u)--(\v);
      }
    }

    \only<2-10>{
      \draw[thArrow,draw=ThGold!80!ThInk,dashed,line width=1.15pt]
        (pa.south) .. controls (3,-2.05) and (8,-2.05) .. (pb.south);
    }
    \only<11>{
      \draw[thArrow,draw=ThGold!80!ThInk,line width=2.0pt]
        (pa.south) .. controls (3,-2.05) and (8,-2.05) .. (pb.south);
    }
    \only<12-17>{
      \draw[thArrowThin,draw=ThMuted!45,dashed,line width=0.80pt]
        (pa.south) .. controls (3,-2.05) and (8,-2.05) .. (pb.south);
    }

    \foreach \step/\u/\v in {
      4/pnext/pbrone,6/pbrone/pbrtwo,7/pbrtwo/pb,8/pb/purone,
      10/pultwo/pa,12/pb/purone,13/purone/purtwo,14/purtwo/pr,
      15/pr/pulone,16/pulone/pultwo,17/pultwo/pa}{
      \only<\step>{
        \draw[thArrow,draw=ThBlue,line width=2.0pt] (\u)--(\v);
      }
    }
  \end{scope}

  \only<2-10>{
    \node[thLabel,text=ThGold!80!ThInk] at (5.5,-1.95)
      {unused auxiliary chord};
  }
  \only<11>{
    \node[thLabel,text=ThGold!80!ThInk] at (5.5,-1.95)
      {used auxiliary chord};
  }
  \only<12-17>{
    \node[thLabel,text=ThMuted!70] at (5.5,-1.95)
      {used auxiliary chord};
  }

  % An inactive probe leaves the pointer in place; otherwise mark the destination.
  \foreach \step/\destination in {
    2/pa,3/pnext,4/pbrone,5/pbrone,6/pbrtwo,7/pb,8/purone,9/pultwo,10/pa,
    11/pb,12/purone,13/purtwo,14/pr,15/pulone,16/pultwo,17/pa}{
    \only<\step>{
      \node[draw=ThGold!85!ThInk,line width=1.6pt,rounded corners=1.8pt,
        inner sep=3pt,fit=(\destination)] {};
    }
  }
\end{tikzpicture}}

\newcommand{\ThImbalanceToMultiplicity}{%
\begin{tikzpicture}[x=1cm,y=1cm]
  \path[use as bounding box] (0,-1.15) rectangle (11.50,1.35);
  \foreach \j in {0,1,2,9,10,11,12}{
    \node[thDet,minimum width=0.40cm,minimum height=0.40cm,inner sep=0pt]
      (s\j) at ({0.45+0.88*\j},0.08) {};
  }
  \foreach \j in {3,...,8}{
    \node[thDet,fill=ThTeal!14,minimum width=0.40cm,
      minimum height=0.40cm,inner sep=0pt]
      (s\j) at ({0.45+0.88*\j},0.08) {};
  }
  \begin{scope}[on background layer]
    % Six consecutive vertices form one visibly extended interval.
    \path[draw=ThTeal!35,fill=ThTeal!7,rounded corners=3pt,line width=0.65pt]
      (2.65,-0.40) rectangle (7.93,0.90);
    \foreach \j [evaluate=\j as \k using int(\j+1)] in {0,...,11}{
      \draw[thArrowThin,draw=ThInkSoft!55,line width=0.65pt]
        (s\j.east)--(s\k.west);
    }
    % Two entering chords above the backbone and one leaving chord below.
    \foreach \u/\v/\h in {0/6/1.25,1/4/0.70}{
      \draw[thArrow,draw=ThTeal,line width=1.15pt]
        (s\u.north) .. controls +(0,\h) and +(0,\h) .. (s\v.north);
    }
    \draw[thArrow,draw=ThCoral,line width=1.15pt]
      (s5.south) .. controls +(0,-1.10) and +(0,-1.10) .. (s12.south);
    % Highlight the two boundary backbone edges whose use counts are compared.
    \draw[thArrowThin,draw=ThGold!85!ThInk,line width=1.35pt]
      (s2.east)--(s3.west);
    \draw[thArrowThin,draw=ThGold!85!ThInk,line width=1.35pt]
      (s8.east)--(s9.west);
  \end{scope}
  \node[thLabel,text=ThTeal!75!ThInk,font=\small] at (6.65,0.65) {\(I\)};
  \node[thLabel,text=ThGold!85!ThInk,font=\small] at (2.65,-0.62)
    {\(c_{\mathrm{in}}\)};
  \node[thLabel,text=ThGold!85!ThInk,font=\small] at (7.93,0.55)
    {\(c_{\mathrm{out}}\)};
\end{tikzpicture}}

% --------------------------------------------------------------------------
% Pairwise objective and finite-memory DP.
% --------------------------------------------------------------------------
\newcommand{\ThPairwiseShortcut}{%
\begin{tikzpicture}[x=1cm,y=1cm]
  \node[thActive] (vi) at (0.65,0) {$v_i$};
  \node[thInactive] (a) at (2.05,0) {};
  \node[thInactive] (b) at (3.20,0) {};
  \node[text=ThMuted] at (4.25,0) {$\cdots$};
  \node[thInactive] (c) at (5.30,0) {};
  \node[thActive] (vj) at (6.80,0) {$v_j$};
  \draw[thFaint] (vi)--(a)--(b);
  \draw[thFaint] (b)--(3.85,0);
  \draw[thFaint] (4.65,0)--(c)--(vj);
  \draw[thArrow,draw=ThBlue] ([yshift=0.28cm]vi.north east)
    to[bend left=22] node[above=0.05cm,thLabel,text=ThBlue]
    {\(p_i p_j\prod_{i<k<j}(1-p_k)\,d(v_i,v_j)\)}
    ([yshift=0.28cm]vj.north west);
  \node[thSmallLabel] at (3.70,-0.65) {the pair contributes only when every interior vertex is inactive};
\end{tikzpicture}}

\newcommand{\ThDPBlocks}{%
\begin{tikzpicture}[x=1cm,y=1cm]
  \path[use as bounding box] (0,-1.00) rectangle (11.60,1.10);
  \tikzset{dpblock/.style={rounded corners=1.5pt,
    minimum width=1.65cm,minimum height=1.08cm,inner sep=0pt,
    line width=0.8pt}}
  % A schematic run of rare-vertex blocks: dots are the vertices in each block.
  \foreach \j/\x in {1/0.95,2/2.95}{
    \node[dpblock,draw=ThMuted!65,fill=ThPaper] (B\j) at (\x,0.40) {};
    \node[thLabel] at (\x,0.65) {$B_{\j}$};
    \foreach \dx in {-0.42,0,0.42}{
      \draw[draw=ThMuted!70,fill=ThMuted!25,line width=0.6pt]
        ({\x+\dx},0.16) circle[radius=0.075cm];
    }
  }
  \node[thLabel] (dots) at (4.45,0.40) {$\cdots$};
  \foreach \j/\x/\lab in {3/5.85/{t-2},4/8.00/{t-1},5/10.15/t}{
    \node[dpblock,draw=ThTeal,fill=ThTeal!8] (B\j) at (\x,0.40) {};
    \node[thLabel,text=ThTeal!65!ThInk] at (\x,0.65) {$B_{\lab}$};
    \foreach \dx in {-0.42,0,0.42}{
      \draw[draw=ThTeal!75!ThInk,fill=ThTeal!35,line width=0.6pt]
        ({\x+\dx},0.16) circle[radius=0.075cm];
    }
  }
  \foreach \u/\v in {B1/B2,B2/dots,dots/B3,B3/B4,B4/B5}{
    \draw[thArrowThin,draw=ThLine] (\u.east)--(\v.west);
  }
  \draw[thArrowThin,draw=ThTeal] (B5.east)--(11.55,0.40);
  \draw[decorate,decoration={brace,amplitude=5pt,mirror},ThMuted,line width=0.8pt]
    (0.125,-0.35)--(3.775,-0.35)
    node[midway,below=0.20cm,thLabel] {earlier blocks};
  \draw[decorate,decoration={brace,amplitude=5pt,mirror},ThTeal,line width=1.0pt]
    (5.025,-0.35)--(10.975,-0.35)
    node[midway,below=0.20cm,thLabel,text=ThTeal] {relevant suffix \(S\)};
\end{tikzpicture}}

% --------------------------------------------------------------------------
% Four-vertex asymmetric clairvoyance-gap example.
% --------------------------------------------------------------------------
\newcommand{\ThAsymFourVertexSeed}{%
\begin{tikzpicture}[x=1cm,y=1cm]
  \path[use as bounding box] (-2.55,-2.02) rectangle (2.55,2.22);
  \node[thDepot,minimum width=0.78cm,minimum height=0.58cm,
    font=\fontsize{9.5pt}{11pt}\selectfont] (r) at (0,1.55) {$r$};
  \node[thDet,minimum size=0.58cm,
    font=\fontsize{9.5pt}{11pt}\selectfont] (a) at (-1.82,0) {$a$};
  \node[thDet,minimum size=0.58cm,
    font=\fontsize{9.5pt}{11pt}\selectfont] (b) at (1.82,0) {$b$};
  \node[thStoch,minimum size=0.60cm,
    font=\fontsize{9.5pt}{11pt}\selectfont] (c) at (0,-1.55) {$c$};

  \begin{scope}[on background layer]
    \draw[thArrowThin,draw=ThInkSoft,line width=0.82pt]
      (r) to[bend right=11] (a);
    \draw[thArrowThin,draw=ThInkSoft,line width=0.82pt]
      (a) to[bend right=11] (r);
    \draw[thArrowThin,draw=ThInkSoft,line width=0.82pt]
      (r) to[bend left=11] (b);
    \draw[thArrowThin,draw=ThInkSoft,line width=0.82pt]
      (b) to[bend left=11] (r);
    \draw[thArrowThin,draw=ThInkSoft,line width=0.82pt] (a)--(b);
    \draw[thArrowThin,draw=ThInkSoft,line width=0.82pt] (b)--(c);
    \draw[thArrowThin,draw=ThInkSoft,line width=0.82pt] (c)--(a);
  \end{scope}

  \node[thSmallLabel,text=ThTeal,anchor=west] at (0.38,-1.55)
    {$p_c=\tfrac12$};
\end{tikzpicture}}

\newcommand{\ThAsymFourVertexOutcomes}{%
\begin{tikzpicture}[x=1cm,y=1cm]
  \path[use as bounding box] (0,-2.20) rectangle (11,1.50);
  \draw[ThLine,line width=0.65pt] (5.50,-2.02)--(5.50,1.43);
  \node[thDepot,minimum width=0.74cm] (lr) at (2.62,1.05) {$r$};
  \node[thDet] (la) at (1.16,-0.18) {$a$};
  \node[thDet] (lb) at (4.08,-0.18) {$b$};
  \node[thInactive] (lc) at (2.62,-1.38) {$c$};
  \begin{scope}[on background layer]
    \foreach \x/\y/\bend in {lr/la/right,la/lr/right,lr/lb/left,lb/lr/left}{
      \draw[thArrowThin,draw=ThLine,line width=0.58pt]
        (\x) to[bend \bend=10] (\y);
    }
    \foreach \x/\y in {la/lb,lb/lc,lc/la}{
      \draw[thArrowThin,draw=ThLine,line width=0.58pt] (\x)--(\y);
    }
    \draw[thArrow,draw=ThTeal,line width=1.48pt]
      (lr) to[bend right=10] (la);
    \draw[thArrow,draw=ThTeal,line width=1.48pt] (la)--(lb);
    \draw[thArrow,draw=ThTeal,line width=1.48pt]
      (lb) to[bend left=10] (lr);
  \end{scope}
  \node[font=\fontsize{9.8pt}{11.5pt}\selectfont,text=ThTeal]
    at (2.62,-1.93) {$r\to a\to b\to r$, cost $3$};

  \node[thDepot,minimum width=0.74cm] (rr) at (8.38,1.05) {$r$};
  \node[thDet] (ra) at (6.92,-0.18) {$a$};
  \node[thDet] (rb) at (9.84,-0.18) {$b$};
  \node[thActive] (rc) at (8.38,-1.38) {$c$};
  \begin{scope}[on background layer]
    \foreach \x/\y/\bend in {rr/ra/right,ra/rr/right,rr/rb/left,rb/rr/left}{
      \draw[thArrowThin,draw=ThLine,line width=0.58pt]
        (\x) to[bend \bend=10] (\y);
    }
    \foreach \x/\y in {ra/rb,rb/rc,rc/ra}{
      \draw[thArrowThin,draw=ThLine,line width=0.58pt] (\x)--(\y);
    }
    \draw[thArrow,draw=ThTeal,line width=1.48pt]
      (rr) to[bend left=10] (rb);
    \draw[thArrow,draw=ThTeal,line width=1.48pt] (rb)--(rc);
    \draw[thArrow,draw=ThTeal,line width=1.48pt] (rc)--(ra);
    \draw[thArrow,draw=ThTeal,line width=1.48pt]
      (ra) to[bend right=10] (rr);
  \end{scope}
  \node[font=\fontsize{9.8pt}{11.5pt}\selectfont,text=ThTeal]
    at (8.38,-1.93) {$r\to b\to c\to a\to r$, cost $4$};
\end{tikzpicture}}

\setbeameroption{hide notes}

\begin{document}

% ===========================================================================
% Title
% ===========================================================================
\begin{frame}[plain]
  \begin{tikzpicture}[remember picture,overlay]
    \fill[ThPaper] (current page.south west) rectangle (current page.north east);

    \node[anchor=west,text=ThInk,text width=11.35cm,align=left,
      font=\fontsize{24pt}{28pt}\selectfont\bfseries]
      at ([xshift=0.98cm,yshift=0.70cm]current page.west)
      {Adaptivity Bounds for the\\
       Stochastic Traveling\\
       Salesperson Problem};
    \node[anchor=south west,text=ThInk,
      font=\fontsize{13.2pt}{16pt}\selectfont\bfseries]
      at ([xshift=0.98cm,yshift=1.08cm]current page.south west)
      {Teng Liu};
    \node[anchor=south west,text=ThMuted,
      font=\fontsize{9.8pt}{12pt}\selectfont]
      at ([xshift=0.98cm,yshift=0.45cm]current page.south west)
      {Advisors: Dr.\ Sharat Ibrahimpur and Prof.\ Dr.\ Vera Traub};
  \end{tikzpicture}
\end{frame}

% ===========================================================================
% Problem and benchmarks
% ===========================================================================
\section{Problem and benchmarks}

\begin{frame}{The traveling salesperson problem}
  \begin{columns}[c,onlytextwidth]
    \begin{column}{0.55\textwidth}
      \begin{itemize}
        \setlength{\itemsep}{0.62em}
        \item We are given a finite vertex set \(V\), a root \(r\), and
          metric travel costs \(d\).
        \item A tour starts at \(r\), visits every other vertex, and returns
          to \(r\).
        \item The goal is to minimize the total travel cost.
      \end{itemize}
    \end{column}
    \begin{column}{0.40\textwidth}
      \centering
      \resizebox{0.94\linewidth}{!}{\ThTSPGraph}
    \end{column}
  \end{columns}
  \thsource{Thesis, Chapter 1}
\end{frame}

\begin{frame}{Stochastic TSP}
  \begin{columns}[c,onlytextwidth]
    \begin{column}{0.57\textwidth}
      \begin{itemize}
        \setlength{\itemsep}{0.46em}
        \item Each non-root vertex \(v\) is independently active with
          probability \(p_v\in[0,1]\); the root \(r\) is always active.
        \item Only active vertices require service.
        \item Start and end at \(r\), visit every active vertex, and minimize
          expected travel cost.
        \item If every \(p_v=1\), this is the classical TSP.
      \end{itemize}
    \end{column}
    \begin{column}{0.38\textwidth}
      \centering
      \resizebox{0.96\linewidth}{!}{\ThStochasticGraph}
    \end{column}
  \end{columns}
  \thsource{Thesis, Chapter 1}
\end{frame}

\begin{frame}{Three benchmarks}
  \begin{itemize}
    \setlength{\itemsep}{0.52em}
    \item \textbf{A posteriori --- full information:} observe the complete
      active set, then choose the best tour.
    \item \textbf{A priori --- no activation information:} fix one order
      before observing activations, then skip inactive vertices.
    \item \textbf{Adaptive --- information revealed over time:} probe one
      vertex at a time; if the probed vertex is active, visit it immediately;
      later probes may depend on earlier outcomes.
  \end{itemize}
  \thsource{Thesis, Chapter 1}
\end{frame}

\begin{frame}{Notation}
  \begin{columns}[T,onlytextwidth]
    \begin{column}{0.48\textwidth}
      \textbf{Instance and realization}
      \begin{itemize}\setlength{\itemsep}{0.48em}
        \item \(\mathcal I\): a stochastic TSP instance.
        \item \(V\): the vertex set; \(r\): the always-active root.
        \item \(V^\circ:=V\setminus\{r\}\), \(n:=|V^\circ|\): the
          non-root vertices and their number.
        \item \(d(u,v)\): the travel cost from \(u\) to \(v\).
        \item \(p_v\in[0,1]\): the activation probability of
          \(v\in V^\circ\).
      \end{itemize}
    \end{column}
    \begin{column}{0.48\textwidth}
      \textbf{Active sets and optimal costs}
      \begin{itemize}\setlength{\itemsep}{0.48em}
        \item \(A\sim\mathcal Q\): the realized active set and its
          distribution.
        \item \(\OPT(A)\): the minimum tour cost on \(A\cup\{r\}\).
        \item \(\optpost\): optimal a posteriori expected cost.
        \item \(\optadapt\): optimal adaptive expected cost.
        \item \(\optprior\): optimal a priori expected cost.
      \end{itemize}
    \end{column}
  \end{columns}
  \thsource{Thesis, Chapters 1--2}
\end{frame}

\begin{frame}{A posteriori}
  \centering
  \resizebox{0.82\linewidth}{!}{\ThPosterioriExample}

  \medskip
  For each active set \(A\), choose a minimum-cost tour on
  \(A\cup\{r\}\).
  \[
    \optpost=\E_{A\sim\mathcal Q}\!\left[\OPT(A)\right].
  \]
  \thsource{Thesis, Chapter 1}
\end{frame}

\begin{frame}{A priori}
  \centering
  \resizebox{0.82\linewidth}{!}{\ThPrioriExample}

  \medskip
  Fix \(T\) before seeing \(A\); then delete the inactive vertices.
  \[
    \optprior=\min_T\E_{A\sim\mathcal Q}
      \!\left[\cost(T\vert_A)\right].
  \]
  \thsource{Thesis, Chapter 1}
\end{frame}

\begin{frame}{Adaptive}
  \centering
  \resizebox{0.82\linewidth}{!}{\ThProbeRealizations}

  \medskip
  A policy \(\pi\) chooses the next probe from earlier outcomes and has
  expected cost \(\cost(\pi)\).
  \[
    \optadapt=\min_\pi\cost(\pi).
  \]
  \thsource{Thesis, Section 2.1}
\end{frame}

\begin{frame}{Three gaps}
  \centering
  \resizebox{0.78\linewidth}{!}{\ThGapArcs}

  \vspace{-0.05cm}
  \begin{itemize}\setlength{\itemsep}{0.40em}
    \item \textbf{\textcolor{ThTeal}{Clairvoyance gap:}}
      \(\displaystyle
        \operatorname{CG}(\mathcal I)
        :=\frac{\optadapt(\mathcal I)}{\optpost(\mathcal I)}.\)
    \item \textbf{\textcolor{ThBlue}{Adaptivity gap:}}
      \(\displaystyle
        \operatorname{AG}(\mathcal I)
        :=\frac{\optprior(\mathcal I)}{\optadapt(\mathcal I)}.\)
    \item \textbf{\textcolor{ThGold!80!ThInk}{Obliviousness gap:}}
      \(\displaystyle
        \operatorname{OG}(\mathcal I)
        :=\frac{\optprior(\mathcal I)}{\optpost(\mathcal I)}.\)
  \end{itemize}
  \thsource{Thesis, Section 2.2}
\end{frame}

\begin{frame}{Asymmetric \(\optpost\) v.s.\ \(\optadapt\)}
  \begin{columns}[c,onlytextwidth]
    \begin{column}{0.50\textwidth}
      \centering
      \resizebox{0.82\linewidth}{!}{\ThAsymFourVertexSeed}
    \end{column}
    \begin{column}{0.46\textwidth}
      \begin{itemize}\setlength{\itemsep}{0.64em}
        \item The root \(r=a\) and vertices \(b,c\) are always active.
        \item Vertex \(d\) is active independently with probability \(1/2\).
        \item The arc \(b\to c\) remains direct, while its reverse is
          replaced by \(c\to d\to b\).
      \end{itemize}
    \end{column}
  \end{columns}
  \thsource{Thesis, Section 5.1}
\end{frame}

\begin{frame}{Asymmetric \(\optpost\) v.s.\ \(\optadapt\)}
  \centering
  \resizebox{0.76\linewidth}{!}{\ThAsymFourVertexOutcomes}

  \vspace{0.10cm}
  \small
  An adaptive policy chooses its first probe before learning whether \(d\) is
  active, so it cannot realize both highlighted optima. Probing \(c,d,b\) costs
  \(4\) in either realization; the fixed order \(a,c,d,b,a\) does as well.
  \[
    \optpost=\tfrac12\cdot3+\tfrac12\cdot4=\frac72,\qquad
    \optadapt=\optprior=4,\qquad
    \operatorname{CG}=\frac{4}{7/2}=\frac87.
  \]
  \thsource{Thesis, Section 5.1}
\end{frame}

\begin{frame}{Symmetric \(\optprior\) v.s.\ \(\optadapt\)}
  \begin{columns}[T,onlytextwidth]
    \begin{column}{0.49\textwidth}
      \centering
      \resizebox{0.90\linewidth}{!}{\ThKFourSubdivision}
    \end{column}
    \begin{column}{0.47\textwidth}
      \vspace{0.38cm}
      \begin{itemize}\setlength{\itemsep}{0.58em}
        \item Take a  \(K_4\) on \(r,a,b,c\), and subdivide
          the edge \(ra\) by \(s\).
        \item The four original vertices are always active; only subdivision vertex \(s\) is
          active with \(p_s=1/2\).
      \end{itemize}
    \end{column}
  \end{columns}
\end{frame}

\begin{frame}{Symmetric \(\optprior\) v.s.\ \(\optadapt\)}
  \centering
  \resizebox{0.90\linewidth}{!}{\ThKFourOutcomes}

  \vspace{0.05cm}
  % \begin{columns}[T,onlytextwidth]
  %   \begin{column}{0.58\textwidth}
      \small Probe \(s\) at the root. If it is inactive, remain at \(r\);
      if it is active, visit \(s\) immediately.
    % \end{column}
    % \begin{column}{0.38\textwidth}
      \[
        \optadapt=\tfrac12\cdot4+\tfrac12\cdot5=\frac92.
      \]
  %   \end{column}
  % \end{columns}
\end{frame}

\begin{frame}{Symmetric \(\optprior\) v.s.\ \(\optadapt\)}
  % \begin{columns}[T,onlytextwidth]
    % \begin{column}{0.56\textwidth}
      % \vspace{0.30cm}
      \begin{itemize}\setlength{\itemsep}{0.78em}
        \item A priori must fix one cyclic order before its activation is known.
             Wherever \(s\) is placed in that fixed order, one realization
          is suboptimal: the same order cannot achieve costs \(4\) and \(5\).
        \item If in some order, $s$ is placed as $\cdots, r, s, a, \cdots$, then
            when $s$ is inactive, we pay cost $2$ for one move from $r$ to $a$; otherwise
            $s$ must be placed adjacent to some vertex with distance at least $2$, thus
            when $s$ is active, we pay cost $2$. $\optprior=5$.
        \item $\text{Adaptivity Gap}=\frac{\optprior}{\optadapt}
        =\color{ThBlue}\frac{10}{9}.$
      \end{itemize}
    % \end{column}
    % \begin{column}{0.40\textwidth}
    %   \[
    %     \text{Adaptivity Gap}=\frac{\optprior}{\optadapt}
    %     =\color{ThBlue}\frac{10}{9}.
    %   \]
    % \end{column}
  % \end{columns}
  \medskip{}
  \thresult[ThBlue]{Even one uncertain vertex can make a priori solution suffers.}
\end{frame}

\begin{frame}{A priori algorithms v.s.\ \(\optpost\) and \(\optprior\)}
  \centering
  Let \(C_{\mathrm{ALG}}\) denote the expected cost of the returned
  a priori tour.

  \vspace{0.10cm}
  {\fontsize{10.8pt}{12.7pt}\selectfont
  \renewcommand{\arraystretch}{1.16}
  \setlength{\tabcolsep}{7pt}
  \begin{tabular}{@{}>{\raggedright\arraybackslash}p{4.05cm}
                      >{\centering\arraybackslash}p{2.45cm}
                      >{\raggedright\arraybackslash}p{5.25cm}@{}}
    \toprule
     & Compared with & Guarantee \\
    \midrule
    Shmoys--Talwar (2008) \hfill symmetric
      & \(\optpost\)
      & \(C_{\mathrm{ALG}}/\optpost\le4\) randomized;
        \(\le8\) deterministic \\
    \cmidrule(lr){1-3}
    van Zuylen (2011) \hfill symmetric
      & \(\optprior\)
      & \(C_{\mathrm{ALG}}/\optprior\le6.5\) deterministic \\
    \cmidrule(lr){1-3}
    Blauth et al. (2024) \hfill symmetric
      & \(\optprior\)
      & \(C_{\mathrm{ALG}}/\optprior<3.1\) randomized;
        \(<5.9\) deterministic \\
    \cmidrule(lr){1-3}
    Christalla et al. (2025) \hfill asymmetric
      & \(\optprior\)
      & \(C_{\mathrm{ALG}}/\optprior=O(\log^8 n)\), randomized
        quasipolynomial time \\
    \bottomrule
  \end{tabular}}

\end{frame}

\begin{frame}{Prior information-gap bounds}
  \ThGapProgressTable{0}
  \thsource{Thesis, Sections 1.1 and 1.3}
\end{frame}

% ===========================================================================
% Adaptivity gap
% ===========================================================================
\section{Adaptivity gap}

\begin{frame}{Symmetric adaptivity gap}
  \ThGapProgressTable{1}
  \thsource{Thesis, Theorem 3.7; Shmoys--Talwar}
\end{frame}

\begin{frame}{Symmetric adaptivity gap: random-metric toy example}
  \centering
  \resizebox{0.86\linewidth}{!}{\ThRandomMetricToy}

  \medskip
   Teal edges have distance \(1\); dashed edges have distance \(2\). 
\end{frame}

\begin{frame}{Symmetric adaptivity gap: random-metric construction}
  \begin{columns}[c,onlytextwidth]
    \begin{column}{0.36\textwidth}
      \centering
      \textbf{Random short-edge graph}\par
      \resizebox{0.78\linewidth}{!}{\ThRandomMetric}
    \end{column}
    \begin{column}{0.60\textwidth}
      \begin{itemize}\setlength{\itemsep}{0.48em}
        \item Let \(V=\{r\}\cup[n]\) and sample
          \(G\sim G(n+1,\theta)\): every unordered pair is independently sampled
          with probability \(\theta=(\log n)^{-1}\). Call the sampled pair \emph{short-edge}.
        \item A sampled pair has distance \(1\), and other pairs have
          distance \(2\). 
        \item Independently of the metric, activate each non-root vertex
          with probability \(p=(\log n)^{-4}\).
      \end{itemize}
    \end{column}
  \end{columns}

\end{frame}

\begin{frame}{Symmetric adaptivity gap: a posteriori benchmark}
      \begin{itemize}\setlength{\itemsep}{0.52em}
        \item Conditional on \(|A|=k\), the short-edge graph on
          \(A\cup\{r\}\) is exactly \(G(k+1,\theta)\).
        \item Typically \(K\approx np\), and
          \[
            \theta K\asymp\frac{n}{(\log n)^5}\gg\log K.
          \]

         With high probability, the short-edge graph on \(A\) is Hamiltonian.
      \end{itemize}

  \[
    \E_D[\optpost(D)]=(1+o(1))W.
  \]
\end{frame}

\begin{frame}{Symmetric adaptivity gap: a priori benchmark}
    \resizebox{0.75\linewidth}{!}{\ThRandomShortcutPair}

    \medskip
    Intuitively, consider a fixed a priori tour containing some pair $(u,w)$, as shown above, the distance $d(u,v) = 2$ with high probability.

    Summing these activation-weighted pair costs over a fixed a priori tour \(T\)
    gives
    \[
        \E_D[\text{cost}_D(T)]=(2-\theta)W.
    \]

    Weighted Hoeffding shows that one fixed tour is lucky
    with probability \(\ll1/n!\).

    Union bound over all orders therefore gives \(\optprior(D)\ge(2-o(1))W\) with high probability.
\end{frame}

\begin{frame}{Symmetric adaptivity gap: adaptive benchmark}
  \begin{columns}[c,onlytextwidth]
    \begin{column}{0.47\textwidth}
      \centering
      \resizebox{0.96\linewidth}{!}{\ThShortNeighborPolicy}
    \end{column}
    \begin{column}{0.49\textwidth}
      \begin{itemize}\setlength{\itemsep}{0.58em}
        \item At the current active vertex, probe every unprobed
          distance-\(1\) neighbor first.
        \item Inactive: remain in place. Active: visit immediately.
        \item Use a distance-\(2\) move only when no short move is found.
      \end{itemize}
    \end{column}
  \end{columns}

  \medskip
  The analysis shows that, with high probability over the sampled metric,
  \[
    \optadapt(D)\le(1+o(1))W.
  \]
  The policy does not know the active Hamiltonian cycle; it constructs a
  nearly all-short route from the outcomes seen so far.
  \thsource{Thesis, Lemmas 3.12--3.14}
\end{frame}

\begin{frame}{Symmetric adaptivity gap: factor-two lower bound}
  \centering
  \resizebox{0.76\linewidth}{!}{\ThGoodMetricIntersection}

  \medskip
  \[
    \frac{\optprior(d_n)}{\optadapt(d_n)}
    \ge \frac{(2-o(1))W}{(1+o(1))W}=2-o(1).
  \]
  \thsource{Thesis, Theorem 3.7}
\end{frame}

\begin{frame}{Symmetric adaptivity gap: sampled master tour}
  \begin{columns}[c,onlytextwidth]
    \begin{column}{0.43\textwidth}
      \centering
      \resizebox{0.96\linewidth}{!}{\ThSampledBackbone}
    \end{column}
    \begin{column}{0.53\textwidth}
      \small
      Sample backbone \(R\) by including \(r\) definitely and each non-root vertex \(u\) with
      probability \(p_u\).

      \medskip
      Let \(T_R\) be an optimal tour on \(R\), and define
      \[
        D_R(u):=\min_{w\in R\setminus\{u\}}d(u,w).
      \]

      Add two copies of one nearest attachment for every
      \(u\in V^\circ\setminus R\). An Euler tour followed by shortcutting
      gives a fixed a priori tour \(M_R\).
    \end{column}
  \end{columns}
\end{frame}

\begin{frame}{Symmetric adaptivity gap: sampled master tour}
  \begin{columns}[c,onlytextwidth]
    \begin{column}{0.43\textwidth}
      \centering
      \resizebox{0.96\linewidth}{!}{\ThSampledBackbone}
    \end{column}
    \begin{column}{0.53\textwidth}
        \begin{align*}
            C_R &:=\E_A[\cost(M_R(A))\mid R] \\
                &\le \cost(T_R)
          +2\!\sum_{u\in V^\circ\setminus R}p_uD_R(u).
      \end{align*}

      % Symmetry lets the same attachment carry the trip out and back.

    \end{column}
  \end{columns}
\end{frame}

\begin{frame}{Symmetric adaptivity gap: indicator setup}
  For each \(u\in V^\circ\), define
  \[
    I_u:=\mathbf 1_{\{u\in R\}},
    \qquad
    \E_R[I_u]=p_u.
  \]

  \medskip
  The distance
  \[
    D_R(u)=\min_{w\in R\setminus\{u\}}d(u,w)
  \]
  depends only on \(R\setminus\{u\}\), so \(I_u\) and \(D_R(u)\) are
  independent.
\end{frame}

\begin{frame}{Symmetric adaptivity gap: one-vertex comparison}
    From independence we can see
  \[
  \begin{aligned}
    \E_R[p_u(1-I_u)D_R(u)]
      &=p_u(1-p_u)\E_R[D_R(u)],
      \\[0.18cm]
    \E_R[I_uD_R(u)]
      &=p_u\E_R[D_R(u)]
  \end{aligned}
  \]

  Since \(1-p_u\le1\), the first term is no larger than the
  second:
  \[
    \E_R[p_u(1-I_u)D_R(u)]
      \le \E_R[I_uD_R(u)].
  \]
\end{frame}

\begin{frame}{Symmetric adaptivity gap: summing the charges}
  Sum the one-vertex inequality over all \(u\in V^\circ\):
  \[
  \begin{aligned}
    \E_R\!\left[
      \sum_{u\in V^\circ\setminus R}p_uD_R(u)
    \right]
      &=\sum_{u\in V^\circ}
        \E_R[p_u(1-I_u)D_R(u)]\\
      &\le\sum_{u\in V^\circ}\E_R[I_uD_R(u)]\\
      &=\E_R\!\left[
        \sum_{u\in R\setminus\{r\}}D_R(u)
      \right].
  \end{aligned}
  \]

  Here \(1-I_u\) selects \(V^\circ\setminus R\), while \(I_u\) selects
  \(R\setminus\{r\}\).
  \thsource{Shmoys--Talwar, Section 2.2; Blauth et al., Section 1.1}
\end{frame}

\begin{frame}{Symmetric adaptivity gap: tour-edge charging}
  Fix \(R\), and let \(u^+\) be the successor of \(u\) on \(T_R\). Since
  \(u^+\in R\setminus\{u\}\), the definition of \(D_R(u)\) gives
  \[
    D_R(u)\le d(u,u^+).
  \]

  Summing over all sampled non-root vertices,
  \[
    \sum_{u\in R\setminus\{r\}}D_R(u)
      \le\sum_{u\in R\setminus\{r\}}d(u,u^+)
      \le\cost(T_R).
  \]

  Combining this with the previous slide and taking expectations,
  \[
    \E_R\!\left[
      \sum_{u\in V^\circ\setminus R}p_uD_R(u)
    \right]
    \le \E_R[\cost(T_R)]
    =\optpost.
  \]

  \thresult[ThBlue]{The doubled detours contribute at most \(2\optpost\).}
  \thsource{Shmoys--Talwar, Section 2.2; Blauth et al., Section 1.1}
\end{frame}

\begin{frame}{Symmetric adaptivity gap: factor-three upper bound}
  The expected cost over the auxiliary sample is
  \[
    \E_R[C_R]
    \le
    \underbrace{\E_R[\cost(T_R)]}_{\text{backbone }=\optpost}
    +2\underbrace{\E_R\!\left[
      \sum_{u\in V^\circ\setminus R}p_uD_R(u)
    \right]}_{\text{attachments }\le\optpost}
    \le3\optpost.
  \]

  \medskip
  Therefore
    \[
        \optprior \le 3 \optpost
    \]
  \smallskip
\end{frame}

\begin{frame}{Asymmetric adaptivity gap}
  \ThGapProgressTable{2}
  \thsource{Thesis, Theorem 7.1; Christalla--Puhlmann--Traub}
\end{frame}

\begin{frame}{Asymmetric adaptivity gap: construction}
    \makebox[\linewidth][c]{
      \resizebox{0.69\linewidth}{!}{\ThRectGrid}
    }

    \(k^2\) columns and \(k\) vertices in each column, with \(p=1/k, n=k^3\).
\end{frame}

\begin{frame}{Asymmetric adaptivity gap: adaptive benchmark}
  \makebox[\linewidth][c]{%
    \resizebox{0.60\linewidth}{!}{\ThColumnSweep}%
  }

  \vspace{0.12cm}
  Each sweep visits at most one active stochastic vertex per column.

  \medskip
  Let \(X_j\sim\operatorname{Binomial}(k,1/k)\) count the active vertices in
  column \(j\), and let \(M=\max_jX_j\). Then
  \[
    \cost(\pi,A)\le2k^2M,
    \qquad \E[M]=O(\log k).
  \]
Therefore \(\optadapt=O(k^2\log k)\).
\end{frame}

\begin{frame}{Asymmetric adaptivity gap: represented columns}
  \small
  Split the fixed a priori order into \(k^2\) packets of size \(k\).
  Focus on \(P=(z_1,\ldots,z_k)\).

  \vspace{0.04cm}
  \makebox[\linewidth][c]{%
    \resizebox{0.82\linewidth}{!}{\ThPacketColumns}%
  }

  \vspace{0.30cm}

  Choose one representative from each represented column in packet order:
  \(y_1,\ldots,y_s\).
  \thsource{Adapted from Thesis, Section 7.2}
\end{frame}

\begin{frame}{Asymmetric adaptivity gap: packet distance}
  \begin{columns}[T,onlytextwidth]
    \begin{column}{0.485\textwidth}
      {\fontsize{11.2pt}{13.2pt}\selectfont\bfseries\color{ThBlue}
        \(s\ge k/2\): many represented columns\par}

      \vspace{0.04cm}
      \makebox[\linewidth][c]{%
        \resizebox{0.90\linewidth}{!}{\ThManyRepresentedCase}%
      }

      {\fontsize{10.2pt}{12.2pt}\selectfont
        For fixed \(a\), these are \(s-a\) distinct positive integers:
        \[
          \begin{aligned}
            \sum_{b>a}d_k(y_a,y_b)
            &\ge1+\cdots+(s-a),\\
            \sum_{a<b}d_k(z_a,z_b)
            &\ge\binom{s+1}{3}\ge\frac{k^3}{96}.
          \end{aligned}
        \]
      }
    \end{column}
    \begin{column}{0.485\textwidth}
      {\fontsize{11.2pt}{13.2pt}\selectfont\bfseries\color{ThCoral}
        \(s<k/2\): repeated columns\par}

      \vspace{0.04cm}
      \makebox[\linewidth][c]{%
        \resizebox{0.90\linewidth}{!}{\ThRepeatedColumnsCase}%
      }

      {\fontsize{10.2pt}{12.2pt}\selectfont
        At least \(k-s>k/2\) packet pairs share a column.
        Each has distance \(k^2\):
        \[
          \sum_{a<b}d_k(z_a,z_b)
          >\frac{k}{2}\,k^2.
        \]
      }
    \end{column}
  \end{columns}

\end{frame}

\begin{frame}{Asymmetric adaptivity gap: a priori lower bound}
    \vspace{0.30cm}
  \makebox[\linewidth][c]{%
    \resizebox{0.65\linewidth}{!}{\ThPacketShortcutEvent}%
  }

  For \(a<b\) in one packet, the shortcut \(z_a\to z_b\) is realised when
  both endpoints are active and every intervening vertex is inactive.
  \[
    \Pr[z_a\!\to z_b\text{ used}]
    =p^2(1-p)^{b-a-1}
    \ge\frac{1}{4k^2},
    \qquad p=\frac1k.
  \]
  \vspace{0.18cm}
  \[
    \E[\text{one packet's contribution}]
    \ge\frac{1}{4k^2}\cdot\frac{k^3}{96}=\frac{k}{384},
    \;
    \optprior\ge k^2\cdot\frac{k}{384}=\frac{k^3}{384}.
  \]
  % \vspace{-0.16cm}
  %   \[
  %     \operatorname{AG}(n)
  %     \ge\operatorname{AG}(k^3)
  %     \ge\frac{\optprior}{\optadapt}
  %     =\Omega\!\left(\frac{k}{\log k}\right)
  %     =\Omega\!\left(\frac{n^{1/3}}{\log n}\right).
  %   \]
\end{frame}

\begin{frame}{Asymmetric adaptivity gap: ratio}
  Combining the two bounds, with \(n=k^3\), gives
  \[
    \operatorname{AG}(n)
    \ge \frac{\optprior}{\optadapt}
    =\Omega\!\left(\frac{k^3}{k^2\log k}\right)
    =\Omega\!\left(\frac{k}{\log k}\right)
    =\Omega\!\left(\frac{n^{1/3}}{\log n}\right).
  \]
  \thsource{Adapted from Thesis, Theorem 7.1}
\end{frame}

% ===========================================================================
% Clairvoyance gap
% ===========================================================================
\section{Clairvoyance gap}

\begin{frame}{Symmetric clairvoyance gap}
  \ThGapProgressTable{3}
  \thsource{Thesis, Theorem 4.6; Shmoys--Talwar}
\end{frame}

\begin{frame}{Symmetric clairvoyance gap: construction}
  \begin{columns}[T,onlytextwidth]
    \begin{column}{0.62\textwidth}
      \centering
      \resizebox{0.96\linewidth}{!}{\ThThreeChain}
    \end{column}
    \begin{column}{0.35\textwidth}
      \vspace{0.32cm}
      Three always-active chains \(A,B,C\), each of length \(H=2m\).

      \vspace{0.26cm}
      The only semi-active vertex is \(s\), with \(p_s=1/2\).

      \vspace{0.26cm}
      Cross-chain lengths:
      \[
        U=3m+2,\qquad P=4m.
      \]
    \end{column}
  \end{columns}
  \thsource{Thesis, Section 4.1}
\end{frame}

\begin{frame}{Symmetric clairvoyance gap: realizations}
  \centering
  \resizebox{0.92\linewidth}{!}{\ThOrderConflict}

  \vspace{0.32cm}
  \thresult[ThBlue]{The posterior solution sees the realization and chooses the
  cheaper order; an adaptive policy has a common prefix before probing \(s\).}
  \thsource{Thesis, Section 4.1}
\end{frame}

\begin{frame}{Symmetric clairvoyance gap: lower bound}
  \begin{columns}[T,onlytextwidth]
    \begin{column}{0.51\textwidth}
      \begin{thtightitemize}
        \item Before probing \(s\), both executions traverse the same edges.
        \item Any sufficiently cheap inactive execution needs the \(A,B,C\)
          crossing pattern.
        \item Any sufficiently cheap active execution needs the incompatible
          \(A,C,s,B\) pattern.
      \end{thtightitemize}
    \end{column}
    \begin{column}{0.45\textwidth}
      \vspace{0.05cm}
      \small
      \[
        \begin{aligned}
          \optpost  &=18m+4,\\
          \optadapt &=20m+4.
        \end{aligned}
      \]
      \[
        \frac{\optadapt}{\optpost}
        =\frac{20m+4}{18m+4}
        \longrightarrow
        \color{ThTeal}\frac{10}{9}.
      \]
    \end{column}
  \end{columns}

\end{frame}

\begin{frame}{Asymmetric clairvoyance gap}
  \ThGapProgressTable{4}
  \thsource{Thesis, Theorem 5.5 and Corollary 6.2}
\end{frame}

\begin{frame}{Asymmetric clairvoyance gap: lower bound}
  \begin{columns}[c,onlytextwidth]
    \begin{column}{0.54\textwidth}
      \centering
      \resizebox{0.98\linewidth}{!}{\ThLayeredStage}
    \end{column}
    \begin{column}{0.42\textwidth}
      \raggedright
      {\bfseries\color{ThTeal}A posteriori}\par
      \smallskip
      \small
      Reveal the realization first, then choose two compatible traversals.
      \[
        \textcolor{ThTeal}{\le 2+o(1)}
      \]

      \medskip
      {\bfseries\color{ThCoral}Adaptive}\par
      \smallskip
      The policy chooses too early. A bad pair forces a third wraparound.
      \[
        \textcolor{ThCoral}{\ge 3-o(1)}
      \]
    \end{column}
  \end{columns}
\end{frame}

% ===========================================================================
% Asymmetric upper bound
% ===========================================================================
\section{An asymmetric adaptive policy}

\begin{frame}{Asymmetric clairvoyance gap: upper-bound challenge}
  \centering
  \resizebox{0.81\linewidth}{!}{\ThDetourComparison}

  \vspace{0.10cm}
  \begin{columns}[T,onlytextwidth]
    \begin{column}{0.47\textwidth}
      \small In a symmetric metric, a round-trip attachment costs at most twice
      the nearest-backbone distance.
    \end{column}
    \begin{column}{0.47\textwidth}
      \small In an asymmetric cyclic grid, returning to the same anchor can cost
      a full revolution per active outside vertex.
    \end{column}
  \end{columns}
  \thsource{Thesis, Chapter 6 introduction}
\end{frame}

\begin{frame}{Asymmetric clairvoyance gap: sampled backbone}
  \begin{columns}[c,onlytextwidth]
    \begin{column}{0.50\textwidth}
      \centering
      \resizebox{0.82\linewidth}{!}{\ThSampledBackboneCycle}
    \end{column}
    \begin{column}{0.46\textwidth}
      \small
      Fix \(W\subseteq V^\circ\), the non-root vertices of this subinstance.
      Independently of the actual active set, sample
      \[
        S\subseteq W,\qquad \Pp(v\in S)=p_v.
      \]
      Vertices with \(p_v=0\) may first be deleted. Set
      \[
        R=S\cup\{r\}\quad\text{and}\quad U=W\setminus S.
      \]

      An optimal tour \(T_R\) on \(R\) is the sampled backbone. Its cyclic
      order defines the backbone edges and cyclic intervals.
    \end{column}
  \end{columns}

  \centering
  \[
    \optpost(W):=\E[\OPT(A\cap W)]
    =\E_S[\cost(T_R)].
  \]
  \thsource{Thesis, Section 6.1}
\end{frame}

\begin{frame}{Asymmetric clairvoyance gap: circulation}
  \begin{columns}[c,onlytextwidth]
    \begin{column}{0.46\textwidth}
      \centering
      \resizebox{0.82\linewidth}{!}{\ThCirculationRelaxation}
    \end{column}
    \begin{column}{0.50\textwidth}
      \footnotesize
      A detour \(a\to i\to b\) avoids returning to \(a\), but choosing
      \(a,b\) separately can leave surplus flow at one backbone vertex and a
      deficit at another.

      \smallskip
      A closed traversal needs equal incoming and outgoing flow at every
      vertex. Thus, for fixed \(S\), seek \(z^S\ge0\):
      \[
        \begin{aligned}
          z^S(\delta^+(v))&=z^S(\delta^-(v)) &&\forall v,\\[0.14cm]
          z^S(\delta^+(i))&=z^S(\delta^-(i))=p_i
            &&\forall i\in U,\\[0.14cm]
          z^S_{ij}&=0 &&\forall i,j\in U.
        \end{aligned}
      \]
    \end{column}
  \end{columns}

  \vspace{0.10cm}
  \begin{center}
    \small Equal incoming and outgoing flow is exactly the circulation
    condition. The \(p_i\) units through \(i\) pair into detours
    \(a\to i\to b\); after collapse to \(a\to b\), the backbone remains
    balanced.
  \end{center}
  \thsource{Thesis, Section 6.1}
\end{frame}

\begin{frame}{Asymmetric clairvoyance gap: feasibility}
  \begin{columns}[c,onlytextwidth]
    \begin{column}{0.47\textwidth}
      \centering
      \resizebox{0.89\linewidth}{!}{\ThLongestRunWitness}
    \end{column}
    \begin{column}{0.09\textwidth}
      \centering
      \ThTransformArrow
    \end{column}
    \begin{column}{0.36\textwidth}
      \centering
      \resizebox{0.89\linewidth}{!}{\ThParallelDetours}
    \end{column}
  \end{columns}

  {\normalfont\footnotesize
    Sample \(\widetilde A\), set \(C=S\cup\widetilde A\), and cut an optimal
    \(T_C\) at backbone vertices. For a segment with \(k_j\) outside vertices,
    use one detour through each and \(M-k_j\) direct copies, where
    \(M:=\max_j k_j\). Every segment now carries \(M\) units.
  }

  \[
    \underbrace{1\text{ in}=1\text{ out}}_{u_{j,\ell}\in\widetilde A\setminus S}
    \qquad
    \underbrace{M\text{ in}=M\text{ out}}_{s_j\in R}
    \quad\Longrightarrow\quad
    z^{\widetilde A,S}\text{ is a circulation.}
  \]

  {\normalfont\footnotesize
    No \(U\to U\) flow arcs occur. Averaging gives
    \(z^S(\delta^\pm(i))=p_i\), so \(z^S\) is feasible for
    \(\mathrm{CIRC}(S)\).
  }
  \thsource{Thesis, Section 6.1}
\end{frame}

\begin{frame}{Asymmetric clairvoyance gap: circulation cost}
  \begin{columns}[c,onlytextwidth]
    \begin{column}{0.51\textwidth}
      \centering
      \resizebox{0.96\linewidth}{!}{\ThSegmentPadding}
    \end{column}
    \begin{column}{0.45\textwidth}
      \small
      By the triangle inequality, every detour and every direct copy for
      \(P_j\) costs at most the length of \(P_j\). Thus \(P_j\) contributes
      at most \(M\) times its original length, and
      \[
        \cost(z^{\widetilde A,S})\le M\,\OPT(C).
      \]
    \end{column}
  \end{columns}

  \smallskip
  \small Conditioned on \(C\), every vertex independently lies outside the
  backbone with probability at most \(1/2\). Therefore
  \[
    \Pp(M\ge t\mid C)\le n2^{-t},
    \qquad
    \E[M\mid C]=O(\log n).
  \]
  Averaging the cost and using \(\E[\OPT(C)]\le 2\optpost(W)\) gives
  \[
    \E_S[\cost(z^S)]=O(\log n)\,\optpost(W).
  \]
  \thsource{Thesis, Theorem 6.5}
\end{frame}

\begin{frame}{Asymmetric clairvoyance gap: joint rounding}
  \centering
  \small Couple \(z^S\) into detour weights \(x_{i,ab}\), and let
  \(y_{ab}:=z^S_{ab}\) be its backbone-only flow. Then
  \(\sum_{a,b}x_{i,ab}=p_i\) for every outside vertex \(i\).

  \resizebox{0.78\linewidth}{!}{\ThJointRounding}

  \begin{itemize}\small\setlength{\itemsep}{0.42em}
    \item Jointly select one endpoint pair \((a_i,b_i)\) for each outside
      vertex and round \(y_{ab}\) to integral multiplicities \(N_{ab}\) of
      auxiliary chords.
    \item Everything is fixed before the actual activations. Rounding the two
      families separately could create large interval imbalance.
  \end{itemize}
  \thsource{Thesis, Sections 6.2--6.3}
\end{frame}

\begin{frame}{Asymmetric clairvoyance gap: rounding guarantee}
  \centering
  \resizebox{0.60\linewidth}{!}{\ThIntervalImbalance}

  {\footnotesize
    For a cyclic interval \(I\), let \(\chi_I(a,b)=+1\) when the chord
    \(a\to b\) leaves \(I\), and \(-1\) when it enters. Define
    \[
      \Delta_\rho(I):=
      \sum_{i\in U}p_i\chi_I(a_i,b_i)
      +\sum_{a,b\in R}N_{ab}\chi_I(a,b).
    \]
    Joint rounding guarantees, simultaneously,
    \[
      |\Delta_\rho(I)|=O(\log^2 n)\quad\forall I
      \quad\text{in every outcome},
    \]
    \[
      \E_\rho[\text{activation-weighted rounded cost}]=\cost(z^S).
    \]
  }
  \thsource{Thesis, Lemma 6.6}
\end{frame}

\begin{frame}{Asymmetric clairvoyance gap: rounding argument}
  \centering
  \resizebox{0.70\linewidth}{!}{\ThDyadicRounding}

  \begin{itemize}\small\setlength{\itemsep}{0.30em}
    \item Keep every outside-vertex equation. Retain a dyadic interval
      equation only while at least \(L=\Theta(\log n)\) touching variables
      remain unfixed.
    \item Fewer constraints than variables leave a feasible direction. An
      unbiased move to a boundary point preserves expectations and fixes at
      least one more variable.
  \end{itemize}

  {\small
    \[
      \text{dyadic error }O(\log n)
      \ \times\ O(\log n)\text{ pieces}
      \ =\ O(\log^2 n).
    \]
  }
  \thsource{Thesis, Lemma 6.6}
\end{frame}

\begin{frame}{Asymmetric clairvoyance gap: adaptive policy}
  \centering
  \resizebox{0.91\linewidth}{!}{\ThAdaptiveChordWalk}

  \vspace{0.08cm}
  \begin{itemize}\setlength{\itemsep}{0.44em}
    \item At a backbone vertex \(a\), probe the outside vertices assigned
      tail \(a\), one at a time.
    \item After that list is exhausted, the virtual pointer takes an unused
      auxiliary chord leaving \(a\), if one remains; otherwise it takes the
      next backbone edge.
  \end{itemize}
  \centering
  The virtual walk treats \(R\) as active. On its first encounter with a
  sampled backbone vertex, probe it; if it is inactive, the physical route
  shortcuts it.
  \thsource{Thesis, Section 6.4}
\end{frame}

\begin{frame}{Asymmetric clairvoyance gap: activation noise}
  \centering
  \resizebox{0.66\linewidth}{!}{\ThIntervalImbalance}

  {\footnotesize
    Let \(W_\tau=\{v:p_v<\tau\}\), \(U=W_\tau\setminus S\), and
    \(\mu=\sum_{i\in W_\tau}p_i\). Let \(F(A)\) contain every auxiliary chord
    and the assigned chord \(a_i\to b_i\) for each active \(i\in A\cap U\).
    For every cyclic interval \(I\),
    \[
      \sum_{f\in F(A)}\chi_I(f)
      =\underbrace{\Delta_\rho(I)}_{O(\log^2 n)}
      +\underbrace{\sum_{i\in U}
        (\mathbf 1_{\{i\in A\}}-p_i)\chi_I(a_i,b_i)}_{\text{activation noise}}.
    \]
    A maximal partial-sum bound controls all cyclic intervals at once:
    \[
      \E_A\max_I
        \left|\sum_{i\in U}(\mathbf 1_{\{i\in A\}}-p_i)
          \chi_I(a_i,b_i)\right|
      =O(\sqrt\mu).
    \]
  }
  \thsource{Thesis, Lemma 6.9}
\end{frame}

\begin{frame}{Asymmetric clairvoyance gap: policy cost}
  \centering
  \resizebox{0.58\linewidth}{!}{\ThImbalanceToMultiplicity}

  {\footnotesize
    For the two boundary edges \(e,e'\) of \(I\), flow balance gives
    \[
      |c_e-c_{e'}|
      =\left|\sum_{f\in F(A)}\chi_I(f)\right|,
      \qquad
      \Lambda(F(A)):=1+\max_I\left|\sum_{f\in F(A)}\chi_I(f)\right|.
    \]
    The virtual walk uses every chord once and every backbone edge at most
    \(\Lambda(F(A))\) times. After averaging over activations, \(S\), and the
    rounding outcome \(\rho\), the two cost contributions \emph{add}:
    \[
      \begin{aligned}
        \text{detours and auxiliary chords}
          &\lesssim \log n\,\optpost(W_\tau),\\
        \text{repeated backbone edges}
          &\lesssim (\log^2 n+\sqrt\mu)\optpost(W_\tau).
      \end{aligned}
    \]
    Averaging therefore fixes some \((S,\rho)\) with
    \(\cost(\pi_\tau^{S,\rho})
      =O(\!\log^2 n+\sqrt\mu\!)\optpost(W_\tau)\).
  }
  \thsource{Thesis, Lemmas 6.8--6.9}
\end{frame}

\begin{frame}{Asymmetric clairvoyance gap: upper bound}
  \begin{columns}[T,onlytextwidth]
    \begin{column}{0.48\textwidth}
      \small
      \theyebrow{rare vertices}
      \[
        \begin{aligned}
          W_\tau&=\{v:p_v<\tau\},\\
          \mu&=\sum_{v\in W_\tau}p_v\le n\tau.
        \end{aligned}
      \]
      The circulation policy costs
      \[
        O\!\left(\log^2 n+\sqrt{n\tau}\right)\optpost(W_\tau).
      \]
    \end{column}
    \begin{column}{0.48\textwidth}
      \small
      \theyebrow{high-probability vertices}
      \[
        H=\{v:p_v\ge\tau\}.
      \]
      A fixed-tour policy costs
      \[
        O(1/\tau)\optpost(H).
      \]
    \end{column}
  \end{columns}

  \vspace{0.12cm}
  \centering
  Concatenate the two rooted policies and choose \(\tau=n^{-1/3}\):
  \[
    \cost(\pi)=O(n^{1/3})\optpost,
    \qquad
    \operatorname{CG}(n)=O(n^{1/3}).
  \]
  \smallskip
  \thpill[ThCoral]{existence result; no polynomial-time claim}
  \thsource{Thesis, Theorem 6.1}
\end{frame}

% ===========================================================================
% A priori algorithm
% ===========================================================================
\section{Almost-optimal a priori tours}

\begin{frame}{A priori approximation}
  \centering
  Let \(C_{\mathrm{ALG}}\) denote the expected cost of the returned
  a priori tour.

  \vspace{0.08cm}
  {\fontsize{9.6pt}{11.3pt}\selectfont
  \renewcommand{\arraystretch}{1.08}
  \setlength{\tabcolsep}{6pt}
  \begin{tabular}{@{}>{\raggedright\arraybackslash}p{4.05cm}
                      >{\centering\arraybackslash}p{2.25cm}
                      >{\raggedright\arraybackslash}p{5.45cm}@{}}
    \toprule
     & Compared with & Guarantee \\
    \midrule
    Shmoys--Talwar (2008) \hfill symmetric
      & \(\optpost\)
      & \(C_{\mathrm{ALG}}/\optpost\le4\) randomized;
        \(\le8\) deterministic \\
    \cmidrule(lr){1-3}
    van Zuylen (2011) \hfill symmetric
      & \(\optprior\)
      & \(C_{\mathrm{ALG}}/\optprior\le6.5\) deterministic \\
    \cmidrule(lr){1-3}
    Blauth et al. (2024) \hfill symmetric
      & \(\optprior\)
      & \(C_{\mathrm{ALG}}/\optprior<3.1\) randomized;
        \(<5.9\) deterministic \\
    \cmidrule(lr){1-3}
    Christalla et al. (2025) \hfill asymmetric
      & \(\optprior\)
      & \(C_{\mathrm{ALG}}/\optprior=O(\log^8 n)\), randomized
        quasipolynomial time \\
    \cmidrule(lr){1-3}
    \shortstack[l]{\textcolor{ThResultBlue}{\bfseries This thesis}\\
      symmetric and asymmetric}
      & \textcolor{ThResultBlue}{\(\optprior\)}
      & \thgapbox[minimum width=5.18cm,minimum height=1.02cm,%
          text width=4.88cm,align=left,draw=ThResultBlue,%
          fill=ThResultBlue!10,line width=1.4pt]{\textcolor{ThResultBlue}{%
          \(\boldsymbol{C_{\mathrm{ALG}}/\optprior\le1+\varepsilon}\);\\
          deterministic, \(\boldsymbol{c_\varepsilon^n\poly(n)}\) time}} \\
    \bottomrule
  \end{tabular}}
  \thsource{Thesis, Theorem 8.5}
\end{frame}

\begin{frame}{A priori approximation: compressed dynamic program}
  \centering
  \resizebox{0.58\linewidth}{!}{\ThDPBlocks}

  \begin{itemize}\small\setlength{\itemsep}{0.34em}
    \item Truncate negligible shortcut contributions and group light vertices
      into low-activation-mass blocks.
    \item The dynamic program remembers only the relevant suffix and one
      boundary block.
  \end{itemize}

  \centering
  {\small
    \[
      \cost(\widehat T)\le(1+\varepsilon)\optprior,
      \qquad
      \text{time }c_\varepsilon^n\poly(n).
    \]
    Deterministic for symmetric and asymmetric metrics.
  }
  \thsource{Thesis, Theorem 8.5}
\end{frame}

% ===========================================================================
% Closing
% ===========================================================================
\section{Takeaway}

\begin{frame}{Takeaways}
  \centering
  \theyebrow{worst-case gaps}

  \vspace{0.08cm}
  \renewcommand{\arraystretch}{1.52}
  \setlength{\tabcolsep}{12pt}
  \begin{tabular}{>{\bfseries}lcc}
    \toprule
    Metric & \(\optprior/\optadapt\) & \(\optadapt/\optpost\) \\
    \midrule
    Symmetric
      & \(2\le\operatorname{AG}\le3\)
      & \(10/9\le\operatorname{CG}\le3\) \\
    \cmidrule(lr){1-3}
    Asymmetric
      & \(\displaystyle
          \Omega(\frac{n^{1/3}}{\log n})
          \le \operatorname{AG}
          \le O(\sqrt n)\)
      & \(\displaystyle
          \frac32
          \le \operatorname{CG}
          \le O(n^{1/3})\) \\
    \bottomrule
  \end{tabular}
\end{frame}

\begin{frame}{Open questions: gap bounds}
  \begin{itemize}
    \setlength{\itemsep}{0.82em}
    \item \textbf{Symmetric clairvoyance gap.}
      What is its exact supremum?
    \item \textbf{Symmetric adaptivity gap.}
      Is it exactly \(2\), or can it be larger?
    \item \textbf{Asymmetric adaptivity gap.}
      Can the lower bound \(\Omega(n^{1/3}/\log n)\) or the upper bound
      \(O(\sqrt n)\) be improved?
    \item \textbf{Asymmetric clairvoyance gap.}
      Is it bounded by a constant, or must it grow with~\(n\)?
  \end{itemize}
\end{frame}

\begin{frame}{Open questions: algorithms}
  \begin{itemize}
    \setlength{\itemsep}{1.18em}
    \item \textbf{Exact algorithms for a priori TSP.}
      Can exact optimization beat enumeration of the \(n!\) orders in
      \mbox{\(2^{o(n\log n)}\) time}? More strongly, is there an exact
      \mbox{\(c^n\poly(n)\)-time algorithm} for some constant~\(c\)?
    \item \textbf{Adaptive algorithms for asymmetric stochastic TSP.}
      Can one compute in polynomial time an adaptive policy~\(\pi\) with
      \(\cost(\pi)\le\operatorname{polylog}(n)\cdot\optpost\)?
  \end{itemize}
\end{frame}

% ===========================================================================
% Backup
% ===========================================================================
\appendix
\section{Backup}

\begin{frame}{Backup: Adaptive DP function}
  Let \(J(u,S)\) be the optimal expected remaining cost from current position
  \(u\) with unprobed set \(S\). Then
  \[
    J(u,S)
    =\min_{v\in S}
      \left[
        (1-p_v)J(u,S\setminus\{v\})
        +p_v\bigl(d(u,v)+J(v,S\setminus\{v\})\bigr)
      \right],
  \]
  with \(J(u,\varnothing)=d(u,r)\).

  \vspace{0.35cm}
  The state contains only the current position and the set of unprobed
  vertices. 
\end{frame}

\begin{frame}{Backup: factor-two symmetric adaptivity lower bound}
  \small Fix \(T=(r,v_1,\ldots,v_n,r)\) and write \(q=1-p\). The shortcut
  coefficients are
  \[
    \alpha_i=p\bigl(q^{i-1}+q^{n-i}\bigr)
    \quad\text{for }\{r,v_i\},
    \qquad
    \beta_{ij}=p^2q^{j-i-1}
    \quad\text{for }\{v_i,v_j\}.
  \]
  They satisfy
  \[
    \sum_e c_e(T)=W,
    \qquad
    \sum_e c_e(T)^2\le2np^3.
  \]
  For \(D(e)=2-X_e\) with independent
  \(X_e\sim\operatorname{Bernoulli}(\theta)\), weighted Hoeffding gives
  \[
    \Pp_D\!\left[
      C_D(T)\le\left(2-\frac{2}{\log n}\right)W
    \right]
    \le e^{-n(\log n)^2}.
  \]
  Hence the probability that \emph{some} fixed order is cheaper is at most
  \[
    n!e^{-n(\log n)^2}=o(1).
  \]
  \thsource{Thesis, Lemmas 3.9--3.11}
\end{frame}

\begin{frame}{Backup: the backbone circulation LP}
  For a sampled backbone \(R\) and outside set \(U\), the circulation sends
  \(p_i\) units through every \(i\in U\):
  \[
    \begin{aligned}
      \min\quad &\sum_{a,b\in R\cup U} d(a,b)z_{ab}\\
      \text{s.t.}\quad
      &\sum_b z_{vb}=\sum_a z_{av} &&\forall v,\\
      &\sum_a z_{ai}=\sum_b z_{ib}=p_i &&\forall i\in U,\\
      &z_{ij}=0 &&\forall i,j\in U,\\
      &z_{ab}\ge0 &&\forall a,b.
    \end{aligned}
  \]
  Flow through \(i\) is coupled into a distribution over detours
  \(a\to i\to b\).
  \thsource{Thesis, Section 6.1}
\end{frame}

\begin{frame}{Backup: cyclic-interval discrepancy}
  \small
  Joint rounding selects one pair per outside vertex and preserves expected
  cost.
  \medskip

  For a cyclic interval \(I\) of the backbone tour, define
  \(\chi_I(a,b):=\mathbf 1_{\{a\in I\}}-\mathbf 1_{\{b\in I\}}\).
  \par\smallskip
  It equals \(+1\) when the chord \(a\to b\) leaves \(I\), \(-1\) when it
  enters, and \(0\) when \(a,b\) lie on \mbox{the same side of~\(I\)}.
  \par\smallskip
  Joint rounding guarantees
  \[
    \left|
      \sum_{i\in U}p_i\chi_I(a_i,b_i)
      +\sum_{a,b\in R}N_{ab}\chi_I(a,b)
    \right|
    =O(\log^2 n).
  \]
\end{frame}

\end{document}
